\documentclass[11pt,oneside]{book}
\usepackage{notes-preamble}
\pgfplotsset{compat=1.18,compat/show suggested version=false}

\allowdisplaybreaks
\setcounter{secnumdepth}{2}
\setcounter{tocdepth}{2}
\emergencystretch=3em
\hbadness=10000
\vbadness=10000
\titleformat{\chapter}
  {\fontsize{22}{26}\normalfont\bfseries\raggedright}
  {Chapter \thechapter}
  {22pt}
  {\uline{#1}}

% Standard amsthm environments, numbered consecutively within each section.
\theoremstyle{plain}
\newtheorem{notetheorem}{Theorem}[section]
\newtheorem{noteproposition}[notetheorem]{Proposition}
\newtheorem{notelemma}[notetheorem]{Lemma}
\newtheorem{notecorollary}[notetheorem]{Corollary}

\theoremstyle{definition}
\newtheorem{notedefinition}[notetheorem]{Definition}
\newtheorem{noteexercise}[notetheorem]{Exercise}
\newtheorem{noteexample}[notetheorem]{Example}
\newtheorem{notealgorithm}[notetheorem]{Algorithm}

\theoremstyle{remark}
\newtheorem{noteremark}[notetheorem]{Remark}

% Retain the original color-box presentation around the numbered statements.
\tcolorboxenvironment{notedefinition}{boxrule=0pt,boxsep=0pt,colback=red!10,left=8pt,right=8pt,enhanced jigsaw,borderline west={2pt}{0pt}{red},sharp corners,before skip=10pt,after skip=10pt,breakable}
\tcolorboxenvironment{noteproposition}{boxrule=0pt,boxsep=0pt,colback=Orange!10,left=8pt,right=8pt,enhanced jigsaw,borderline west={2pt}{0pt}{Orange},sharp corners,before skip=10pt,after skip=10pt,breakable}
\tcolorboxenvironment{notetheorem}{boxrule=0pt,boxsep=0pt,colback=blue!10,left=8pt,right=8pt,enhanced jigsaw,borderline west={2pt}{0pt}{blue},sharp corners,before skip=10pt,after skip=10pt,breakable}
\tcolorboxenvironment{notelemma}{boxrule=0pt,boxsep=0pt,colback=Cyan!10,left=8pt,right=8pt,enhanced jigsaw,borderline west={2pt}{0pt}{Cyan},sharp corners,before skip=10pt,after skip=10pt,breakable}
\tcolorboxenvironment{notecorollary}{boxrule=0pt,boxsep=0pt,colback=violet!10,left=8pt,right=8pt,enhanced jigsaw,borderline west={2pt}{0pt}{violet},sharp corners,before skip=10pt,after skip=10pt,breakable}
\tcolorboxenvironment{noteremark}{boxrule=0pt,boxsep=0pt,blanker,borderline west={2pt}{0pt}{Green},left=8pt,right=8pt,before skip=10pt,after skip=10pt,breakable}
\tcolorboxenvironment{noteexample}{boxrule=0pt,boxsep=0pt,blanker,borderline west={2pt}{0pt}{Black},left=8pt,right=8pt,before skip=10pt,after skip=10pt,breakable}
\tcolorboxenvironment{noteexercise}{boxrule=0pt,boxsep=0pt,colback=gray!8,left=8pt,right=8pt,enhanced jigsaw,borderline west={2pt}{0pt}{gray},sharp corners,before skip=10pt,after skip=10pt,breakable}
\tcolorboxenvironment{notealgorithm}{boxrule=0pt,boxsep=0pt,colback=Yellow!12,left=8pt,right=8pt,enhanced jigsaw,borderline west={2pt}{0pt}{Goldenrod},sharp corners,before skip=10pt,after skip=10pt,breakable}


\newcommand{\Dom}{\operatorname{D}}
\newcommand{\Ran}{\operatorname{Ran}}
\newcommand{\Ker}{\operatorname{Ker}}
\newcommand{\Graph}{\operatorname{Graph}}
\newcommand{\Span}{\operatorname{span}}
\newcommand{\dist}{\operatorname{dist}}
\newcommand{\diam}{\operatorname{diam}}
\newcommand{\supp}{\operatorname{supp}}
\newcommand{\Int}{\operatorname{int}}
\newcommand{\conv}{\operatorname{conv}}
\newcommand{\aco}{\operatorname{aco}}
\newcommand{\weakto}{\rightharpoonup}
\newcommand{\weakstarto}{\overset{*}{\rightharpoonup}}
\newcommand{\Bop}[2]{\mathcal{B}(#1,#2)}
\newcommand{\Kop}[2]{\mathcal{K}(#1,#2)}
\newcommand{\Id}{I}
\newcommand{\closure}[1]{\overline{#1}}
\newcommand{\ann}[1]{#1^{\perp}}
\newcommand{\preann}[1]{{}^{\perp}\!#1}
\newcommand{\spec}{\sigma}
\newcommand{\res}{\rho}
\newcommand{\spr}{r}
\newcommand{\dd}{\,\mathrm{d}}
\newcommand{\sign}{\operatorname{sgn}}
\newcommand{\Var}{\operatorname{Var}}
\newcommand{\card}{\operatorname{card}}


\title{\Huge \textbf{Functional Analysis}}
\author{Wei Chen}
\date{16 July 2026}

\begin{document}
\maketitle
\allowdisplaybreaks
	\pagenumbering{roman}
	\setcounter{page}{1}
	\begin{center}
		\Large\textbf{Preface} 

        
        
    \end{center}~\
        This is the lecture note for the course \textit{Functional Analysis} which followed the textbook \textit{Introductory Functional Analysis with Applications} by Erwin Kreyszig. The course is taught by Prof. Yutian Li in 2023 Spring at the Chinese University of Hong Kong, Shenzhen. The note is edited and revised with the support of ChatGPT. The materials presume a basic knowledge of real analysis and general topology, so contents that are close to the boundary of these prerequisites may be omitted. Topics of the note include: normed spaces and Banach spaces, Hilbert spaces, fundamental theorems in functional analysis, and spectral theory.
	\newpage
	\tableofcontents

\mainmatter
\chapter{Metric Spaces}

\section{Isometries and completions}

\begin{notedefinition}[Isometric maps and isometric spaces]
Let $(X,d_X)$ and $(Y,d_Y)$ be metric spaces.  A map $T\colon X\to Y$ is an
\emph{isometry} if
\[
  d_Y(Tx,Ty)=d_X(x,y)
  \qquad(x,y\in X).
\]
The spaces $X$ and $Y$ are \emph{isometric} if there is a bijective isometry
from one onto the other.
\end{notedefinition}

\begin{notetheorem}[Completion]
Every metric space $(X,d)$ has a completion: there is a complete metric space
$\widehat X$ and an isometry $J\colon X\to\widehat X$ whose range is dense.
The completion is unique up to a unique isometry fixing the embedded copy of
$X$.
\end{notetheorem}

\begin{noteexercise}
What is the completion of a discrete metric space?
\end{noteexercise}

\begin{proof}
For the discrete metric,
\[
  d(x,y)=
  \begin{cases}
    0,&x=y,\\
    1,&x\ne y,
  \end{cases}
\]
every Cauchy sequence is eventually constant: taking $\varepsilon=1/2$
forces all sufficiently late terms to be equal.  Hence every Cauchy sequence
already converges in $X$, so $X$ is complete and its completion is $X$ itself.
\end{proof}

\chapter{Normed Spaces and Banach Spaces}

\section{Vector spaces, Hamel bases, and norms}

\begin{notetheorem}[Existence and cardinality of Hamel bases]
Let $X\ne\{0\}$ be a vector space over a field $\mathbb K$.
\begin{enumerate}
  \item $X$ has a Hamel basis.
  \item Any two Hamel bases of $X$ have the same cardinality.
\end{enumerate}
\end{notetheorem}

\begin{proof}
For existence, order the linearly independent subsets of $X$ by inclusion.
The union of a chain is linearly independent, so Zorn's lemma gives a maximal
linearly independent set $E$.  Maximality implies $\Span E=X$.

For uniqueness of cardinality, let $E$ and $F$ be bases.  Every element of
$E$ is a finite linear combination of elements of $F$.  The infinite
Steinitz exchange theorem, equivalently the standard basis-cardinality
argument, gives an injection $E\hookrightarrow F$.  Interchanging $E$ and
$F$ gives an injection $F\hookrightarrow E$.  The Cantor--Bernstein theorem
then yields $|E|=|F|$.
\end{proof}

\begin{noteexercise}[Additive functions on $\mathbb R$]
Let $\mathcal A$ be the set of functions $f\colon\mathbb R\to\mathbb R$
satisfying Cauchy's equation
\[
  f(x+y)=f(x)+f(y).
\]
Describe $\mathcal A$ and determine its cardinality.
\end{noteexercise}

\begin{proof}
Such a function is exactly a $\mathbb Q$-linear map on the
$\mathbb Q$-vector space $\mathbb R$.  Fix a Hamel basis $H$ of
$\mathbb R$ over $\mathbb Q$.  An arbitrary function $a\colon H\to\mathbb R$
extends uniquely to an additive function by
\[
  f\!\left(\sum_{j=1}^n q_jh_j\right)=\sum_{j=1}^nq_ja(h_j).
\]
Conversely every additive function is obtained this way.  Since
$|H|=|\mathbb R|=\mathfrak c$, one has
\[
  |\mathcal A|=|\mathbb R|^{\mathfrak c}=2^{\mathfrak c}.
\]
Only the choices $a(h)=ch$ give continuous additive functions.
\end{proof}

\begin{notelemma}[Extension of a linearly independent family]
Every linearly independent family in a vector space is contained in a Hamel
basis.
\end{notelemma}

\begin{proof}
Apply Zorn's lemma to the linearly independent sets containing the given
family.  A maximal member spans the whole space.
\end{proof}

\begin{notedefinition}[Equivalent norms]
Two norms $\|\cdot\|_1$ and $\|\cdot\|_2$ on a vector space $X$ are
\emph{equivalent} if there are constants $c,C>0$ such that
\[
  c\|x\|_1\le\|x\|_2\le C\|x\|_1
  \qquad(x\in X).
\]
\end{notedefinition}

\begin{notetheorem}[Equivalent norms and topology]
Two norms on $X$ are equivalent if and only if they induce the same topology.
\end{notetheorem}

\begin{proof}
If the topologies agree, the $\|\cdot\|_2$-unit ball is a neighbourhood of
$0$ for $\|\cdot\|_1$, so some $r>0$ satisfies
$B_1(0,r)\subseteq B_2(0,1)$.  Applying this inclusion to
$rx/(2\|x\|_1)$ gives
$\|x\|_2\le 2r^{-1}\|x\|_1$ for $x\ne0$.  Reversing the roles of the
norms gives the other inequality.

Conversely, the displayed two-sided estimate implies
\[
  B_1\!\left(x,\frac rC\right)\subseteq B_2(x,r),
  \qquad
  B_2(x,cr)\subseteq B_1(x,r).
\]
Thus the open sets for either norm are open for the other.
\end{proof}

\begin{notetheorem}[Separable spaces are increasing unions]
Let $X$ be a separable infinite-dimensional normed space.  Then there are
finite-dimensional subspaces
\[
  X_1\subsetneq X_2\subsetneq\cdots,
  \qquad \dim X_n=n,
\]
such that $\bigcup_nX_n$ is dense in $X$.
\end{notetheorem}

\begin{proof}
Let $(d_m)$ be dense.  Scan this sequence in order and retain $d_m$ exactly
when it is not in the span of the previously retained vectors.  Denote the
resulting linearly independent sequence by $(x_n)$ and put
$X_n=\Span\{x_1,\ldots,x_n\}$.  Every discarded $d_m$ already belongs to
one of these spans, so $\bigcup_nX_n$ contains a dense subset.  Since $X$ is
infinite-dimensional, the retained sequence is infinite, and
$\dim X_n=n$.
\end{proof}

\begin{notetheorem}[Every vector space can be normed]
Every vector space admits a norm.
\end{notetheorem}

\begin{proof}
Choose a Hamel basis $(e_i)_{i\in I}$.  Every $x\in X$ has a unique finite
expansion $x=\sum_{i\in F}a_ie_i$.  Define
\[
  \|x\|=\sum_{i\in F}|a_i|.
\]
Uniqueness of the expansion makes this well defined, and the norm axioms are
immediate.
\end{proof}

\section{Distance to compact sets and spaces of functions}

\begin{noteproposition}[Distance to a compact set]
Let $K$ be a nonempty compact subset of a normed space $X$ and define
\[
  d_K(x)=\inf_{z\in K}\|x-z\|.
\]
Then the infimum is attained for every $x$, and
\[
  |d_K(x)-d_K(y)|\le\|x-y\|.
\]
In particular, $d_K$ is continuous.
\end{noteproposition}

\begin{proof}
Choose $z_n\in K$ with $\|x-z_n\|\to d_K(x)$.  Compactness gives a
convergent subsequence $z_{n_k}\to z\in K$, and continuity of the norm gives
$\|x-z\|=d_K(x)$.  For any $z\in K$,
\[
  d_K(x)\le\|x-z\|\le\|x-y\|+\|y-z\|.
\]
Taking the infimum over $z$ gives $d_K(x)-d_K(y)\le\|x-y\|$; interchange
$x$ and $y$.
\end{proof}

\begin{noteremark}[Nearest-point selections]
A nearest point need not be unique, and an arbitrary selection
$P(x)\in K$ with $\|x-P(x)\|=d_K(x)$ need not be continuous.  For example,
for $K=\{-1,1\}\subset\mathbb R$, any nearest-point selection jumps at
$0$.  What is always continuous is the distance function $d_K$.
\end{noteremark}

\begin{notetheorem}[Supremum norm on continuous maps]
Let $K$ be compact and $Y$ a normed space.  The formula
\[
  \|f\|_\infty=\sup_{x\in K}\|f(x)\|
\]
defines a norm on $C(K;Y)$.
\end{notetheorem}

\begin{notetheorem}[Bounded maps]
For arbitrary $X$ and a normed space $Y$, let $B(X;Y)$ be the vector space of
bounded maps $f\colon X\to Y$.  Then $\|f\|_\infty$ is a norm on
$B(X;Y)$.
\end{notetheorem}

\begin{notetheorem}[Uniform limits preserve continuity]
If $f_n\colon X\to Y$ are continuous and $f_n\to f$ uniformly, then $f$ is
continuous.
\end{notetheorem}

\begin{proof}
Fix $x\in X$ and $\varepsilon>0$.  Choose $n$ with
$\|f-f_n\|_\infty<\varepsilon/3$, then choose a neighbourhood $U$ of $x$
so that $\|f_n(y)-f_n(x)\|<\varepsilon/3$ for $y\in U$.  The triangle
inequality gives $\|f(y)-f(x)\|<\varepsilon$.
\end{proof}

\begin{notecorollary}
The subspace $B(X;Y)\cap C(X;Y)$ is closed in $B(X;Y)$ for the supremum
norm.
\end{notecorollary}

\section{Banach spaces and completeness}

\begin{notedefinition}[Banach space]
A normed space is a \emph{Banach space} if it is complete for the metric
$d(x,y)=\|x-y\|$.
\end{notedefinition}

\begin{notedefinition}[Schauder basis]
A sequence $(e_n)$ in a normed space $X$ is a \emph{Schauder basis} if every
$x\in X$ has a unique expansion
\[
  x=\sum_{n=1}^{\infty}a_ne_n
\]
converging in norm.  Unlike a Hamel basis, a Schauder basis uses infinite
series and is therefore a topological notion.
\end{notedefinition}

\begin{noteremark}
A normed space with a Schauder basis is separable, but not every separable
Banach space has a Schauder basis.
\end{noteremark}

\begin{notelemma}[Translation and scalar invariance]
The metric induced by a norm satisfies
\[
  d(x+a,y+a)=d(x,y),
  \qquad
  d(\alpha x,\alpha y)=|\alpha|d(x,y).
\]
\end{notelemma}

\begin{proof}
Both identities follow immediately from
$\|(x+a)-(y+a)\|=\|x-y\|$ and
$\|\alpha(x-y)\|=|\alpha|\|x-y\|$.
\end{proof}

\begin{notetheorem}[Closed subspaces of Banach spaces]
Let $X$ be a Banach space and $Y\subseteq X$ a linear subspace.  Then $Y$ is
Banach with the inherited norm if and only if $Y$ is closed in $X$.
\end{notetheorem}

\begin{proof}
If $Y$ is closed, a Cauchy sequence in $Y$ converges in $X$ and its limit lies
in $Y$.  Conversely, if $Y$ is complete and $y_n\in Y$ converges in $X$ to
$x$, then $(y_n)$ is Cauchy in $Y$ and converges there to some $y\in Y$.
Uniqueness of limits gives $x=y$.
\end{proof}

\begin{notedefinition}[Series in normed spaces]
The series $\sum_{n=1}^{\infty}x_n$ converges if its partial sums converge.
It is \emph{absolutely convergent} if $\sum_n\|x_n\|<\infty$.
\end{notedefinition}

\begin{notetheorem}[Isomorphic images of Banach spaces]
Let $X$ be Banach, let $Y$ be normed, and let
$A\colon X\to Y$ be a bounded bijection with bounded inverse.  Then $Y$ is
Banach.
\end{notetheorem}

\begin{proof}
If $(y_n)$ is Cauchy in $Y$, then
$\|A^{-1}y_n-A^{-1}y_m\|\le\|A^{-1}\|\|y_n-y_m\|$, so
$A^{-1}y_n\to x\in X$.  Hence $y_n\to Ax$.
\end{proof}

\begin{notetheorem}[Closed range from a lower bound]
Let $X$ be Banach, let $Y$ be normed, and let $A\in\Bop{X}{Y}$.  If
\[
  \|Ax\|\ge c\|x\|
  \qquad(x\in X)
\]
for some $c>0$, then $A$ is injective and $\Ran(A)$ is closed in $Y$.
\end{notetheorem}

\begin{proof}
Injectivity is immediate.  If $Ax_n\to y$, then
$\|x_n-x_m\|\le c^{-1}\|Ax_n-Ax_m\|$, so $(x_n)$ is Cauchy and converges
to some $x\in X$.  Continuity gives $Ax=y$.
\end{proof}

\begin{notetheorem}[Neumann series]
Let $X$ be Banach and $A\in\Bop{X}{X}$ with $\|A\|<1$.  Then
\[
  (I-A)^{-1}=\sum_{n=0}^{\infty}A^n
\]
in operator norm, and
\[
  \|(I-A)^{-1}\|\le\frac1{1-\|A\|}.
\]
\end{notetheorem}

\begin{proof}
The series converges absolutely in the Banach algebra $\Bop{X}{X}$.  Its
partial sums $S_N$ satisfy
$(I-A)S_N=S_N(I-A)=I-A^{N+1}$.  Letting $N\to\infty$ gives the inverse.
\end{proof}

\begin{notecorollary}[The invertible group is open]
If $A\in\Bop{X}{X}$ is invertible and
$\|B-A\|<\|A^{-1}\|^{-1}$, then $B$ is invertible and
\[
  \|B^{-1}-A^{-1}\|
  \le
  \frac{\|A^{-1}\|^2\|B-A\|}
       {1-\|A^{-1}\|\|B-A\|}.
\]
Thus inversion is continuous on the group of invertible operators.
\end{notecorollary}

\begin{notetheorem}[A Schauder basis implies separability]
If a normed space has a Schauder basis, then it is separable.
\end{notetheorem}

\begin{proof}
Let $(e_n)$ be a Schauder basis.  The set of finite rational (or Gaussian
rational) linear combinations of the $e_n$ is countable.  Approximation of
the coefficients in a sufficiently long partial sum shows this set is dense.
\end{proof}

\begin{noteexercise}[A power-small operator]
Let $X$ be Banach and suppose $A\in\Bop{X}{X}$ satisfies $\|A^p\|<1$ for
some $p\ge1$.  Show that $I-A$ is invertible.
\end{noteexercise}

\begin{proof}
Since
\[
  I-A^p=(I-A)(I+A+\cdots+A^{p-1}),
\]
the Neumann theorem makes $I-A^p$ invertible.  The factors commute, so
\[
  (I-A)^{-1}=(I+A+\cdots+A^{p-1})(I-A^p)^{-1}.
\]
\end{proof}

\begin{notetheorem}[Absolute convergence characterizes Banach spaces]
A normed space $X$ is Banach if and only if every series satisfying
$\sum_n\|x_n\|<\infty$ converges in $X$.
\end{notetheorem}

\begin{proof}
In a Banach space the partial sums are Cauchy because tails have norm bounded
by the corresponding scalar tail.

Conversely, let $(x_n)$ be Cauchy.  Choose a subsequence $(x_{n_k})$ with
\[
  \|x_{n_{k+1}}-x_{n_k}\|<2^{-k}.
\]
The series
$x_{n_1}+\sum_k(x_{n_{k+1}}-x_{n_k})$ converges by hypothesis, so the
subsequence converges.  A Cauchy sequence with a convergent subsequence
converges to the same limit.
\end{proof}

\begin{notetheorem}[Completion of a normed space]
Every normed space $X$ has a Banach completion $\widehat X$ in which $X$ is
isometrically embedded as a dense subspace.  It is unique up to an isometric
isomorphism fixing $X$.
\end{notetheorem}

\begin{notedefinition}[Seminorm]
A map $p\colon X\to[0,\infty)$ is a seminorm if
\[
  p(\alpha x)=|\alpha|p(x),
  \qquad
  p(x+y)\le p(x)+p(y).
\]
Then $p(0)=0$ and $|p(x)-p(y)|\le p(x-y)$.
\end{notedefinition}

\begin{noteproposition}[Quotient by the null space]
Let $N=\{x:p(x)=0\}$.  Then $N$ is a subspace and
\[
  \|x+N\|_{X/N}=p(x)
\]
is a well-defined norm on $X/N$.
\end{noteproposition}

\begin{proof}
If $x-y\in N$, then
$|p(x)-p(y)|\le p(x-y)=0$, so the formula is independent of the chosen
representative.  Positive definiteness on the quotient and the remaining norm
axioms follow directly.
\end{proof}

\section{Finite-dimensional normed spaces}

\begin{notelemma}[Linear combinations]
If $x_1,\ldots,x_n$ are linearly independent in a normed space, then there is
$c>0$ such that
\[
  \left\|\sum_{j=1}^n a_jx_j\right\|
  \ge c\sum_{j=1}^n|a_j|
  \qquad(a_1,\ldots,a_n\in\mathbb K).
\]
\end{notelemma}

\begin{proof}
On the compact set
$S=\{a\in\mathbb K^n:\sum_j|a_j|=1\}$, the continuous function
$a\mapsto\|\sum_ja_jx_j\|$ has a minimum.  The minimum is positive, because
zero would give a nontrivial linear dependence.  Homogeneity yields the
estimate.
\end{proof}

\begin{notetheorem}[Completeness]
Every finite-dimensional normed space is Banach.  More generally, every
finite-dimensional subspace of a normed space is complete.
\end{notetheorem}

\begin{proof}
Write $y_m=\sum_{j=1}^na_j^{(m)}e_j$ in a basis.  The preceding lemma shows
that a Cauchy sequence $(y_m)$ has each scalar coordinate
$(a_j^{(m)})$ Cauchy.  The coordinates converge in $\mathbb K$, and the
corresponding vector is the limit of $(y_m)$.
\end{proof}

\begin{notecorollary}
Every finite-dimensional subspace of a normed space is closed.
\end{notecorollary}

\begin{notetheorem}[Equivalence of finite-dimensional norms]
Any two norms on a finite-dimensional vector space are equivalent.
\end{notetheorem}

\begin{proof}
Choose a basis and compare each norm with the coordinate $\ell^1$ norm.  The
upper bounds follow from the triangle inequality, and the lower bounds follow
from the linear-combination lemma.
\end{proof}

\begin{notetheorem}[Inequivalent norms in infinite dimension]
Every infinite-dimensional vector space admits two inequivalent norms.
\end{notetheorem}

\begin{proof}
Choose a Hamel basis containing a countable subfamily $(e_n)$.  For a finite
expansion $x=\sum a_ie_i$, set
\[
  \|x\|_1=\sum_i|a_i|,
  \qquad
  \|x\|_2=\sum_iw_i|a_i|,
\]
where $w_n=n$ on the chosen subfamily and $w_i=1$ elsewhere.  Then
$\|e_n\|_1=1$ but $\|e_n\|_2=n$, so no uniform comparison
$\|x\|_2\le C\|x\|_1$ is possible.
\end{proof}

\section{Compactness}

\begin{notedefinition}[Sequential compactness]
A subset $M$ of a metric space is compact if every sequence in $M$ has a
subsequence converging to an element of $M$.  In metric spaces this is
equivalent to the open-cover definition.
\end{notedefinition}

\begin{notelemma}
Every compact subset of a metric space is closed and bounded.
\end{notelemma}

\begin{proof}
If $M$ were unbounded, one could choose $x_n\in M$ with
$d(x_n,x_0)>n$, and no subsequence could converge.  If a sequence in $M$
converges in the ambient space, compactness gives a subsequence converging to
a point of $M$; uniqueness of the limit puts the original limit in $M$.
\end{proof}

\begin{noteexample}[The converse fails]
In $\ell^2$, the set $\{e_n:n\ge1\}$ of standard unit vectors is closed and
bounded, but not compact, because
\[
  \|e_n-e_m\|_2=\sqrt2\qquad(n\ne m).
\]
Thus it has no convergent subsequence.
\end{noteexample}

\begin{notetheorem}[Heine--Borel in finite dimension]
A subset of a finite-dimensional normed space is compact if and only if it is
closed and bounded.
\end{notetheorem}

\begin{proof}
Use a basis to identify the space linearly and homeomorphically with
$\mathbb K^n$.  Boundedness of the vectors implies boundedness of their
coordinate tuples; Bolzano--Weierstrass gives a convergent coordinate
subsequence, and closedness keeps the limit in the set.
\end{proof}

\begin{notelemma}[Riesz's lemma]
Let $Y$ be a proper closed subspace of a normed space $X$.  For every
$0<\theta<1$ there is $z\in X$ such that
\[
  \|z\|=1,
  \qquad
  \|z-y\|>\theta\quad(y\in Y).
\]
\end{notelemma}

\begin{proof}
Choose $x\notin Y$ and put $d=\dist(x,Y)>0$.  Select $y_0\in Y$ with
$\|x-y_0\|<d/\theta$ and set
$z=(x-y_0)/\|x-y_0\|$.  For $y\in Y$,
\[
  \|z-y\|
  =\frac{\|x-(y_0+\|x-y_0\|y)\|}{\|x-y_0\|}
  \ge\frac d{\|x-y_0\|}>\theta.
\]
\end{proof}

\begin{noteremark}
The strict endpoint $\theta=1$ is impossible in the displayed form, because
$y=0$ gives $\|z-y\|=1$.  If a nearest point to $x$ in $Y$ exists, one may
obtain the sharp non-strict statement
$\dist(z,Y)=1$, equivalently $\|z-y\|\ge1$ for every $y\in Y$.
\end{noteremark}

\begin{notetheorem}[Compact unit ball forces finite dimension]
If the closed unit ball of a normed space $X$ is compact, then $X$ is
finite-dimensional.
\end{notetheorem}

\begin{proof}
If $X$ were infinite-dimensional, repeatedly apply Riesz's lemma to
$Y_n=\Span\{x_1,\ldots,x_n\}$ to obtain unit vectors with
$\|x_n-x_m\|>1/2$ for $n\ne m$.  This bounded sequence has no convergent
subsequence, contradicting compactness.
\end{proof}

\begin{noteremark}
A metric space is compact if and only if it is complete and totally bounded.
\end{noteremark}

\begin{notetheorem}[Banach space of continuous functions]
Let $K$ be compact and $Y$ Banach.  Then $C(K;Y)$ is Banach under the
supremum norm.
\end{notetheorem}

\begin{proof}
If $(f_n)$ is Cauchy in $\|\cdot\|_\infty$, then for each $x$ the sequence
$(f_n(x))$ converges in $Y$.  Let $f(x)$ be the pointwise limit.  Passing to
the limit in the uniform Cauchy estimate shows $f_n\to f$ uniformly, and the
uniform-limit theorem makes $f$ continuous.
\end{proof}

\begin{notecorollary}
If $Y$ is Banach, then $B(X;Y)$ is Banach under the supremum norm.
\end{notecorollary}

\begin{notetheorem}[Best approximation in finite-dimensional subspaces]
Let $Y$ be a finite-dimensional subspace of a normed space $X$.  For every
$x\in X$ there is $y_0\in Y$ such that
\[
  \|x-y_0\|=\dist(x,Y).
\]
If, in addition, $X$ is strictly convex, then $y_0$ is unique.
\end{notetheorem}

\begin{proof}
Choose a minimizing sequence in $Y$.  It is bounded, so finite-dimensional
compactness gives a convergent subsequence; continuity of the norm gives a
minimizer.  If $y_1\ne y_2$ are minimizers and $X$ is strictly convex, then
\[
 \left\|x-\frac{y_1+y_2}{2}\right\|
 =\left\|\frac{(x-y_1)+(x-y_2)}2\right\|
 <\dist(x,Y),
\]
a contradiction.
\end{proof}

\begin{noteexample}
The spaces $\ell^p$ are strictly convex for $1<p<\infty$, but not for
$p=1$ or $p=\infty$.  For example, in $\ell^1_2$, let
$Y=\Span\{(1,1)\}$ and $x=(1,0)$.  Then
\[
  \|x-t(1,1)\|_1=|1-t|+|t|=1
  \qquad(0\le t\le1),
\]
so the best approximant is not unique.
\end{noteexample}

\begin{notetheorem}
Every compact subset of an infinite-dimensional normed space has empty
interior.
\end{notetheorem}

\begin{proof}
If a compact set contained a ball $B(x,r)$, then its closed subset
$\overline B(x,r/2)$ would be compact.  Translation and scaling would make
the closed unit ball compact, contradicting the preceding theorem.
\end{proof}

\section{Linear and bounded operators}

\begin{notedefinition}[Inverse operator]
A linear map $T\colon X\to Y$ is invertible if it is bijective.  Its inverse
$T^{-1}\colon Y\to X$ is then linear.  If $S\colon X\to Y$ and
$T\colon Y\to Z$ are bijective, then
\[
  (TS)^{-1}=S^{-1}T^{-1}.
\]
\end{notedefinition}

\begin{notedefinition}[Bounded linear map]
For normed spaces $V,W$ and a linear map $T\colon V\to W$, define
\[
  \|T\|=\sup_{\|v\|\le1}\|Tv\|.
\]
The map is \emph{bounded} if $\|T\|<\infty$.  The vector space of bounded
linear maps is denoted $\Bop{V}{W}$.
\end{notedefinition}

\begin{notetheorem}[Bounded inverse and a lower bound]
Let $A\colon X\to Y$ be a linear map between normed spaces.  The following
are equivalent.
\begin{enumerate}
  \item $A$ is injective and the inverse
        $A^{-1}\colon\Ran(A)\to X$ is bounded.
  \item There is $c>0$ such that
        $\|Ax\|\ge c\|x\|$ for every $x\in X$.
\end{enumerate}
In that case the best lower-bound constant is $c=\|A^{-1}\|^{-1}$.
\end{notetheorem}

\begin{proof}
If $A^{-1}$ is bounded, then
$\|x\|\le\|A^{-1}\|\|Ax\|$.  Conversely, the lower bound gives
injectivity and, for $y=Ax$,
$\|A^{-1}y\|=\|x\|\le c^{-1}\|y\|$.
\end{proof}

\begin{noteexercise}[A functional on $c_0$]
Let $(\alpha_n)\in\ell^1$ and define
\[
  f(x)=\sum_{n=1}^{\infty}\alpha_nx_n,
  \qquad x=(x_n)\in c_0.
\]
Then $f\in c_0^*$ and
\[
  \|f\|=\sum_{n=1}^{\infty}|\alpha_n|.
\]
If infinitely many $\alpha_n$ are nonzero, the norm need not be attained on
the unit sphere of $c_0$.
\end{noteexercise}

\begin{proof}
The estimate $|f(x)|\le\|x\|_\infty\sum|\alpha_n|$ gives the upper bound.
For $N\ge1$, choose $x^{(N)}$ supported on the first $N$ coordinates with
$x_n^{(N)}$ equal to the phase of $\alpha_n$.  Then
$\|x^{(N)}\|_\infty=1$ and
$|f(x^{(N)})|=\sum_{n\le N}|\alpha_n|$, proving equality after
$N\to\infty$.  Equality for one $x\in c_0$ would require
$|x_n|=1$ with the correct phase whenever $\alpha_n\ne0$, which is impossible
when infinitely many coefficients are nonzero.
\end{proof}

\begin{notetheorem}[Operator norm]
The operator norm makes $\Bop{V}{W}$ a normed space.  It satisfies
\[
  \|ST\|\le\|S\|\,\|T\|,
  \qquad
  \|Tv\|\le\|T\|\,\|v\|.
\]
If $W$ is Banach, then $\Bop{V}{W}$ is Banach.
\end{notetheorem}

\begin{proof}
Only completeness needs comment.  If $(T_n)$ is Cauchy in operator norm, then
$(T_nv)$ is Cauchy in $W$ for each $v$.  Define $Tv=\lim_nT_nv$.
Linearity passes to the limit, the uniform operator bound gives boundedness,
and the Cauchy estimate passes to the limit to show $\|T_n-T\|\to0$.
\end{proof}

\begin{notetheorem}[Finite-dimensional domain]
Every linear map from a finite-dimensional normed space into a normed space is
bounded.
\end{notetheorem}

\begin{proof}
For a basis $(e_j)_{j=1}^n$, write $x=\sum a_je_j$.  The
linear-combination lemma bounds $\sum|a_j|$ by a constant multiple of
$\|x\|$, and therefore
\[
  \|Tx\|\le\sum_j|a_j|\|Te_j\|\le C\|x\|.
\]
\end{proof}

\begin{notetheorem}[Characterization by linear functionals]
A normed space is finite-dimensional if and only if every linear functional
on it is continuous.
\end{notetheorem}

\begin{proof}
The forward direction is the preceding theorem.  If $X$ is
infinite-dimensional, choose linearly independent unit vectors $(e_n)$ and
extend them to a Hamel basis.  Define a linear functional by
$\varphi(e_n)=n$ and by zero on the other basis vectors.  Then
$\|e_n\|=1$ but $|\varphi(e_n)|=n$, so $\varphi$ is unbounded.
\end{proof}

\begin{notetheorem}[Continuity is equivalent to boundedness]
A linear map between normed spaces is continuous if and only if it is bounded.
It is enough that it be continuous at one point.
\end{notetheorem}

\begin{proof}
A bounded map is Lipschitz.  Conversely, continuity at $0$ gives
$\delta>0$ such that $\|x\|<\delta$ implies $\|Tx\|<1$.
For $x\ne0$, apply this to $(\delta/2\|x\|)x$ to obtain
$\|Tx\|\le2\delta^{-1}\|x\|$.  Continuity at an arbitrary point translates
to continuity at $0$ by linearity.
\end{proof}

\begin{notecorollary}
For a bounded linear operator $T$, the null space $\Ker(T)$ is closed.
\end{notecorollary}

\begin{notedefinition}[Compact operator]
A linear map $A\colon X\to Y$ is \emph{compact} if it maps every bounded
sequence in $X$ to a sequence with a convergent subsequence in $Y$.
Equivalently, the image of every bounded set is relatively compact.  A compact
linear map is automatically bounded.
\end{notedefinition}

\begin{notetheorem}[Basic compact-operator criteria]
Let $A\colon X\to Y$ be linear.
\begin{enumerate}
  \item Every compact operator is bounded.
  \item $A$ is compact if and only if every bounded sequence $(x_n)$ has a
        subsequence $(x_{n_k})$ for which $(Ax_{n_k})$ converges.
  \item Every bounded finite-rank operator is compact.
\end{enumerate}
\end{notetheorem}

\begin{proof}
If the image of the unit ball were unbounded, one could choose a sequence whose
image has no convergent subsequence.  The sequential criterion is exactly the
definition of relative compactness in metric spaces.  For finite rank, the
image of a bounded set is bounded in a finite-dimensional space, whose closure
is compact.
\end{proof}

\begin{noteexercise}
The inverse of a bounded bijective operator need not be bounded when the
codomain is incomplete.  Let
\[
  T\colon \ell^1\longrightarrow \Ran(T)\subset\ell^1,
  \qquad
  T(x_n)=\left(\frac{x_n}{n}\right),
\]
where the range carries the inherited $\ell^1$ norm.  Then $T$ is a bounded
bijection onto its range, but
$\|T^{-1}e_n/n\|_1=1$ while $\|e_n/n\|_1=1/n$, so $T^{-1}$ is
unbounded.  Accordingly, $\Ran(T)$ is not complete.
\end{noteexercise}

\begin{notetheorem}[Extension from a dense subspace]
Let $U$ be a dense subspace of a normed space $V$, let $W$ be Banach, and let
$S\in\Bop{U}{W}$.  Then there is a unique
$T\in\Bop{V}{W}$ extending $S$, and
\[
  \|T\|=\|S\|.
\]
\end{notetheorem}

\begin{proof}
For $v\in V$, choose $u_n\in U$ with $u_n\to v$ and define
$Tv=\lim_nSu_n$.  The sequence $(Su_n)$ is Cauchy, and completeness of $W$
provides the limit.  The estimate
$\|Su_n-Su_n'\|\le\|S\|\|u_n-u_n'\|$ shows independence of the
approximating sequence.  Passing to limits gives linearity and
$\|Tv\|\le\|S\|\|v\|$; restriction gives the reverse norm inequality.
Uniqueness follows from density and continuity.
\end{proof}

\begin{noteexample}[Completeness of the codomain is essential]
The identity map from $\mathbb Q$ to $\mathbb Q$, with the usual norm and
viewed on the dense subspace $\mathbb Q\subset\mathbb R$, has no continuous
extension $\mathbb R\to\mathbb Q$.
\end{noteexample}

\begin{notecorollary}[Composition estimate]
If $T\in\Bop{U}{V}$ and $S\in\Bop{V}{W}$, then
\[
  \|S\circ T\|\le\|S\|\,\|T\|.
\]
\end{notecorollary}

\begin{noteexercise}[Subspaces and connectedness]
\begin{enumerate}
  \item A linear subspace containing a nonempty open ball is the whole space.
  \item The only subsets of a normed space that are both open and closed are
        $\varnothing$ and the whole space.
  \item The closure of a linear subspace is again a linear subspace.
\end{enumerate}
\end{noteexercise}

\begin{proof}
For (1), translate and scale the ball to show every vector belongs to the
subspace.  For (2), normed spaces are path connected via line segments.
For (3), take convergent sequences from the subspace and pass addition and
scalar multiplication to the limit.
\end{proof}

\begin{notedefinition}[Quotient norm]
Let $U$ be a subspace of a normed space $V$.  On the algebraic quotient set
$V/U$ define
\[
  \|v+U\|=\inf_{u\in U}\|v+u\|=\dist(v,U).
\]
\end{notedefinition}

\begin{notetheorem}[Quotient spaces]
The quotient formula is a norm if and only if $U$ is closed.  Moreover:
\begin{enumerate}
  \item if $V$ is Banach and $U$ is closed, then $V/U$ is Banach;
  \item if $U$ and $V/U$ are Banach, then $V$ is Banach.
\end{enumerate}
\end{notetheorem}

\begin{proof}
The only nonautomatic norm axiom is positive definiteness:
$\|v+U\|=0$ exactly when $v\in\overline U$.  Thus it holds precisely when
$U$ is closed.

For the first completeness statement, from a summable series in $V/U$ choose
representatives with nearly minimal norm; the resulting corrected series in
$V$ converges and projects to the required quotient sum.  For the converse,
let $\sum v_n$ be absolutely convergent in $V$.  Its quotient series
converges, so after subtracting a suitable $v\in V$ its partial sums approach
$U$.  Choose corrections in $U$ with summable errors; completeness of $U$
then gives convergence in $V$.
\end{proof}

\begin{notetheorem}[Equivalent boundedness criteria]
For a linear map $T\colon V\to W$ between normed spaces, the following are
equivalent:
\begin{enumerate}
  \item $T$ is bounded;
  \item $T$ is continuous at some point;
  \item $T$ is uniformly continuous;
  \item $T^{-1}(B_W(0,1))$ contains a neighbourhood of $0$.
\end{enumerate}
\end{notetheorem}

\begin{proof}
Boundedness gives a global Lipschitz estimate.  Continuity at one point
translates to continuity at zero.  If a ball $B_V(0,\delta)$ is mapped into
$B_W(0,1)$, homogeneity gives $\|Tv\|\le2\delta^{-1}\|v\|$.
\end{proof}

\section{Dual spaces}

\begin{notedefinition}[Algebraic and continuous duals]
The algebraic dual of a vector space $X$ is the vector space of all linear
functionals on $X$.  For a normed space $X$, the \emph{dual space} $X^*$ is
the subspace of bounded linear functionals, with norm
\[
  \|f\|=\sup_{\|x\|\le1}|f(x)|.
\]
\end{notedefinition}

\begin{notetheorem}
The dual $X^*$ of every normed space is Banach.
\end{notetheorem}

\begin{proof}
This is the completeness theorem for $\Bop{X}{\mathbb K}$, since the scalar
field is complete.
\end{proof}

\begin{noteexercise}[Codimension-one kernels]
Let $f\ne0$ be a linear functional on a vector space $X$ and choose
$x_0$ with $f(x_0)\ne0$.
\begin{enumerate}
  \item Every $x\in X$ has a unique representation
        $x=\alpha x_0+y$ with $y\in\Ker f$.
  \item If $f_1,f_2\ne0$ have the same kernel, then $f_1$ is a nonzero scalar
        multiple of $f_2$.
\end{enumerate}
\end{noteexercise}

\begin{proof}
Take $\alpha=f(x)/f(x_0)$ and $y=x-\alpha x_0$.  Uniqueness follows by
applying $f$.  For the second claim, use the same decomposition and compare
the values on a vector outside the common kernel.
\end{proof}

\chapter{Fundamental Theorems for Banach Spaces}

\section{Hahn--Banach theory}

\begin{notedefinition}[Linear functional]
A linear functional on a vector space $V$ is a linear map
$\varphi\colon V\to\mathbb K$.
\end{notedefinition}

\begin{notetheorem}[Bounded functionals and closed kernels]
Let $V$ be a normed space and let $0\ne\varphi\colon V\to\mathbb K$ be
linear.  The following are equivalent:
\begin{enumerate}
  \item $\varphi$ is bounded;
  \item $\varphi$ is continuous;
  \item $\Ker\varphi$ is closed.
\end{enumerate}
Since $\varphi\ne0$, its kernel is automatically a proper subspace.  That
properness alone does not imply boundedness.
\end{notetheorem}

\begin{proof}
Boundedness and continuity are equivalent, and continuity makes the inverse
image of $\{0\}$ closed.  Suppose $\Ker\varphi$ is closed.  Choose
$x_0\notin\Ker\varphi$.  Every $x$ has the unique form
$x=y+\alpha x_0$ with $y\in\Ker\varphi$, and closedness gives
$d:=\dist(x_0,\Ker\varphi)>0$.  Hence
\[
  d|\alpha|
  \le\|\alpha x_0-y'\|
  \le\|x\|
\]
for a suitable $y'\in\Ker\varphi$, while
$|\varphi(x)|=|\alpha|\,|\varphi(x_0)|$.  Thus
$|\varphi(x)|\le |\varphi(x_0)|d^{-1}\|x\|$.
\end{proof}
\sourcepages{46--47}

\begin{notetheorem}[Bases and discontinuous functionals]
Every vector space has a Hamel basis.  Consequently, every
infinite-dimensional normed space has a discontinuous linear functional.
\end{notetheorem}

\begin{proof}
The basis theorem follows from Zorn's lemma.  Choose linearly independent unit
vectors $(e_n)$ and extend them to a Hamel basis.  Prescribe
$\varphi(e_n)=n$ and $\varphi=0$ on the remaining basis vectors.  The
resulting algebraic linear functional is unbounded on the unit sphere.
\end{proof}
\sourcepages{47--48}

\begin{notelemma}[One-dimensional extension step]
Let $V$ be a real normed space, let $U\subseteq V$ be a subspace, let
$\psi\in U^*$, and let $h\notin U$.  There is an extension
$\widetilde\psi\colon U+\mathbb Rh\to\mathbb R$ with
$\|\widetilde\psi\|=\|\psi\|$.
\end{notelemma}

\begin{proof}
Any extension has the form
\[
  \widetilde\psi(u+th)=\psi(u)+tc.
\]
The estimate $|\widetilde\psi(u+th)|\le\|\psi\|\|u+th\|$ is equivalent to
choosing $c$ in the interval
\[
\sup_{u\in U}\bigl(-\|\psi\|\|u+h\|-\psi(u)\bigr)
\le c\le
\inf_{v\in U}\bigl(\|\psi\|\|v+h\|-\psi(v)\bigr).
\]
For $u,v\in U$, the left expression at $u$ is at most the right expression at
$v$ because
\[
  \psi(v)-\psi(u)=\psi(v-u)
  \le\|\psi\|\|v-u\|
  \le\|\psi\|\bigl(\|u+h\|+\|v+h\|\bigr).
\]
Thus the interval is nonempty.
\end{proof}
\sourcepage{49}

\begin{notetheorem}[Hahn--Banach theorem]
Let $X$ be a real vector space, let $Y\subseteq X$ be a subspace, and let
$p\colon X\to\mathbb R$ be sublinear.  If $f\colon Y\to\mathbb R$ is linear
and $f(y)\le p(y)$ on $Y$, then $f$ has a linear extension
$F\colon X\to\mathbb R$ satisfying $F(x)\le p(x)$ for all $x\in X$.
\end{notetheorem}

\begin{proof}
Partially order all dominated extensions by extension.  Unions of chains are
again dominated extensions, so Zorn's lemma gives a maximal one.  If its
domain $M$ is not all of $X$, choose $x_0\notin M$.  An extension to
$M+\mathbb Rx_0$ has the form
$F_\alpha(m+tx_0)=F(m)+t\alpha$.  Domination is equivalent to
\[
  \sup_{m\in M}\bigl(F(m)-p(m-x_0)\bigr)
  \le\alpha\le
  \inf_{m\in M}\bigl(p(m+x_0)-F(m)\bigr).
\]
Sublinearity shows the left endpoint does not exceed the right endpoint,
contradicting maximality.  Hence $M=X$.
\end{proof}
\sourcepages{50,53--54}

\begin{notetheorem}[Complex Hahn--Banach theorem]
Let $X$ be a complex vector space, let $Y\subseteq X$, and let
$p\colon X\to[0,\infty)$ be a seminorm.  If $f\colon Y\to\mathbb C$ is
complex linear and $|f(y)|\le p(y)$, then there is a complex-linear extension
$F\colon X\to\mathbb C$ with $|F(x)|\le p(x)$.
\end{notetheorem}

\begin{proof}
Let $u=\operatorname{Re}f$, regarded as a real-linear functional.  Extend
$u$ by the real theorem to a real-linear $U$ with $|U(x)|\le p(x)$.  Define
\[
  F(x)=U(x)-iU(ix).
\]
Then $F$ is complex linear and extends $f$.  Given $x$ with $F(x)\ne0$,
choose $\theta$ so that $e^{-i\theta}F(x)=|F(x)|$.  Since
$F(e^{-i\theta}x)=e^{-i\theta}F(x)$,
\[
  |F(x)|=\operatorname{Re}F(e^{-i\theta}x)=U(e^{-i\theta}x)
  \le p(e^{-i\theta}x)=p(x).
\]
\end{proof}
\sourcepages{50--51,55--56}

\begin{notecorollary}[Norm-preserving extension]
Let $X$ be a real or complex normed space, let $Y\subseteq X$, and let
$f\in Y^*$.  Then $f$ has an extension $F\in X^*$ satisfying
\[
  \|F\|=\|f\|.
\]
\end{notecorollary}

\begin{proof}
Apply Hahn--Banach with the sublinear functional
$p(x)=\|f\|\|x\|$.  The domination yields $\|F\|\le\|f\|$; restriction to
$Y$ yields the reverse inequality.
\end{proof}
\sourcepage{56}

\begin{notetheorem}[Norming functional]
For every nonzero $x_0$ in a normed space $X$, there is $f\in X^*$ such that
\[
  \|f\|=1,
  \qquad
  f(x_0)=\|x_0\|.
\]
Consequently,
\[
  \|x\|=\sup_{0\ne f\in X^*}\frac{|f(x)|}{\|f\|}
       =\sup_{\|f\|\le1}|f(x)|.
\]
\end{notetheorem}

\begin{proof}
On $\Span\{x_0\}$ define
$g(\alpha x_0)=\alpha\|x_0\|$.  It has norm one.  Extend it by
Hahn--Banach.  The norm formula follows from the general estimate
$|f(x)|\le\|f\|\|x\|$ and this norming functional.
\end{proof}
\sourcepage{57}

\begin{notetheorem}[Taylor--Foguel unique-extension criterion]
A normed space $X^*$ is strictly convex if and only if, for every subspace
$Y\subseteq X$ and every $y^*\in Y^*$, there is at most one norm-preserving
extension of $y^*$ to $X$.
\end{notetheorem}

\begin{proof}[Corrected proof outline]
If $F_1\ne F_2$ are norm-preserving extensions of $y^*$, then
$(F_1+F_2)/2$ is also an extension and has norm at least $\|y^*\|$.
Strict convexity gives a contradiction after normalization.

Conversely, suppose $X^*$ is not strictly convex.  Choose distinct
$F_1,F_2\in S_{X^*}$ with $\|(F_1+F_2)/2\|=1$, and put
$H=(F_1+F_2)/2$ and $Y=\Ker(F_1-F_2)$.  Choose $x_n\in S_X$ and unimodular
scalars so that, after replacing $x_n$ by those scalar multiples,
$H(x_n)\to1$.  Since $\operatorname{Re}F_j(x_n)\le1$ and their average
tends to $1$, both $F_1(x_n)$ and $F_2(x_n)$ tend to $1$.  Hence
$(F_1-F_2)(x_n)\to0$.  For the nonzero functional $F_1-F_2$,
\[
  \dist(x,Y)=\frac{|(F_1-F_2)(x)|}{\|F_1-F_2\|},
\]
so there are $y_n\in Y$ with $\|x_n-y_n\|\to0$.  It follows that
$\|H|_Y\|=1$.  But $F_1|_Y=F_2|_Y=H|_Y$, and $F_1,F_2$ are two distinct
norm-preserving extensions of this common restriction.
\end{proof}
\sourcepages{58--59}

\begin{notetheorem}[Separating a point from a closed subspace]
Let $Y$ be a closed subspace of a normed space $X$, let $x_0\notin Y$, and
put $d=\dist(x_0,Y)$.  Then there is $f\in X^*$ such that
\[
  \|f\|=1,
  \qquad
  f|_Y=0,
  \qquad
  f(x_0)=d.
\]
\end{notetheorem}

\begin{proof}
On $Y+\mathbb Kx_0$ define $g(y+\alpha x_0)=\alpha d$.  Since
$d\le\|x_0-y\|$ for every $y\in Y$,
\[
 |g(y+\alpha x_0)|=|\alpha|d
 \le\|y+\alpha x_0\|.
\]
Thus $\|g\|\le1$, while approximating the distance shows $\|g\|=1$.
Hahn--Banach extends $g$ with the same norm.
\end{proof}
\sourcepage{60}

\begin{notecorollary}[Annihilator criterion for density]
For a subspace $Y\subseteq X$,
\[
  \overline Y=X
  \quad\Longleftrightarrow\quad
  \{f\in X^*:f|_Y=0\}=\{0\}.
\]
\end{notecorollary}

\begin{proof}
If $Y$ is dense, continuity forces every functional vanishing on $Y$ to
vanish on $X$.  If $\overline Y\ne X$, separate a point outside
$\overline Y$ from the closed subspace $\overline Y$.
\end{proof}

\begin{notetheorem}[A separable dual forces a separable space]
If $X^*$ is separable, then $X$ is separable.
\end{notetheorem}

\begin{proof}
Choose a dense sequence $(f_n)$ in $S_{X^*}$ and, for each $n$, choose
$x_n\in S_X$ with $|f_n(x_n)|>1/2$.  Let
$M=\overline{\Span\{x_n:n\ge1\}}$.  If $M\ne X$, Hahn--Banach gives
$f\in S_{X^*}$ vanishing on $M$.  Choose $n$ with $\|f-f_n\|<1/2$.
Then $|f_n(x_n)|=|(f_n-f)(x_n)|<1/2$, a contradiction.
\end{proof}
\sourcepages{61--62}

\subsection{Geometric Hahn--Banach theorems}

\begin{notedefinition}[Minkowski functional]
Let $C\subseteq X$ be open, convex, and contain $0$.  Its gauge is
\[
  p_C(x)=\inf\{t>0:x\in tC\}.
\]
Then $p_C$ is sublinear, continuous, and
\[
  C=\{x:p_C(x)<1\}.
\]
\end{notedefinition}

\begin{proof}
Openness and absorption make $p_C$ finite.  Positive homogeneity follows by
rescaling.  If $x\in(r+\varepsilon)C$ and
$y\in(s+\varepsilon)C$, convexity gives
$x+y\in(r+s+2\varepsilon)C$, so
$p_C(x+y)\le p_C(x)+p_C(y)$.  A norm ball contained in $C$ gives
$p_C(x)\le M\|x\|$, hence continuity.  The displayed description of $C$
follows from openness and convexity.
\end{proof}

\begin{notetheorem}[Separation of two open convex sets]
Let $A,B$ be nonempty, disjoint, open, convex subsets of a real normed space
$X$.  Then there are a nonzero $f\in X^*$ and $\gamma\in\mathbb R$ such that
\[
  f(a)\le\gamma\le f(b)
  \qquad(a\in A,\ b\in B).
\]
\end{notetheorem}

\begin{proof}
Put $C=A-B$.  Then $C$ is open and convex and $0\notin C$.  Choose
$c_0\in C$, set $U=C-c_0$, and let $p_U$ be its Minkowski functional.  The
point $x_0=-c_0$ does not belong to $U$, so $p_U(x_0)\ge1$.  On
$\Span\{x_0\}$ define $g(tx_0)=t p_U(x_0)$ and extend it by Hahn--Banach
to a continuous functional $f\le p_U$.  For $c\in C$ one has
$c-c_0\in U$, and therefore
\[
  f(c)=f(c-c_0)+f(c_0)<1-p_U(x_0)\le0.
\]
Thus $f(a)<f(b)$ for every $a\in A$ and $b\in B$.  Consequently
$\sup_A f\le\inf_B f$; choose $\gamma$ between these two numbers.
\end{proof}
\sourcepages{63--65}

\begin{notetheorem}[Strict separation of a compact set]
Let $A$ be closed and convex and let $K$ be compact and convex in a real
normed space, with $A\cap K=\varnothing$.  Then there are
$f\in X^*\setminus\{0\}$, $\gamma\in\mathbb R$, and $\delta>0$ such that
\[
  f(a)\le\gamma-\delta<\gamma+\delta\le f(k)
  \qquad(a\in A,\ k\in K).
\]
\end{notetheorem}

\begin{proof}
Compactness and closedness give $d=\dist(A,K)>0$.  The open convex
neighbourhoods $A+B(0,d/3)$ and $K+B(0,d/3)$ are disjoint.  Separate them by
the preceding theorem.  The two added balls produce a positive gap in the
functional values.
\end{proof}
\sourcepage{65}

\begin{notetheorem}[Nonunique extensions over a strict closed subspace]
Let $Y$ be a proper closed subspace of a normed space $X$ and let
$\ell\in Y^*$.  Besides a norm-preserving extension, there are continuous
extensions with strictly larger norm.
\end{notetheorem}

\begin{proof}
Choose a norm-preserving extension $L\in X^*$.  Since $Y$ is proper and
closed, Hahn--Banach gives a nonzero $g\in X^*$ with $g|_Y=0$.  Then
$L+tg$ extends $\ell$ for every scalar $t$, and
$\|L+tg\|\to\infty$ as $|t|\to\infty$.
\end{proof}
\sourcepage{66}

\begin{noteexercise}[Continuity of sublinear functionals]
If a sublinear functional $p$ on a normed space is continuous at $0$ and
$p(0)=0$, then it is continuous everywhere.
\end{noteexercise}

\begin{proof}
Subadditivity gives
\[
  p(x)-p(y)\le p(x-y),
  \qquad
  p(y)-p(x)\le p(y-x).
\]
Both right-hand sides tend to zero as $x\to y$.
\end{proof}
\sourcepage{67}

\begin{noteexercise}[Kernel contained in a subspace]
Let $\varphi$ be a nonzero linear functional on $V$.  If a subspace
$U\subseteq V$ contains $\Ker\varphi$, then either $U=\Ker\varphi$ or
$U=V$.
\end{noteexercise}

\begin{proof}
If $x\in U\setminus\Ker\varphi$, then every $v\in V$ can be written
$v=(v-\varphi(v)x/\varphi(x))+\varphi(v)x/\varphi(x)$, with the first term
in $\Ker\varphi\subseteq U$ and the second in $U$.
\end{proof}

\begin{noteexercise}[Uniqueness in the one-dimensional extension]
In the one-dimensional Hahn--Banach step, the extension is unique exactly
when the lower and upper endpoints of the admissible interval coincide.
\end{noteexercise}

\begin{noteexercise}[Separating an open ball from the origin]
If an open ball $B(f_0,r)$ in a normed space does not contain $0$, then there
is $\varphi\in X^*$ such that
\[
  \operatorname{Re}\varphi(f)>0
  \qquad(f\in B(f_0,r)).
\]
\end{noteexercise}

\begin{proof}
Since $0\notin B(f_0,r)$, we have $\|f_0\|>r$.  By Hahn--Banach, choose
$\varphi\in X^*$ with $\|\varphi\|=1$ and
$\varphi(f_0)=\|f_0\|$; in the complex case, rotate a norming functional by a
unimodular scalar to obtain this normalization.  If $f=f_0+h$ with
$\|h\|<r$, then
\[
  \operatorname{Re}\varphi(f)
  =\|f_0\|+\operatorname{Re}\varphi(h)
  \ge \|f_0\|-|\varphi(h)|
  >\|f_0\|-r
  >0.
\]
\end{proof}
\sourcepages{68--69}

\begin{noteexercise}[Sublinear level sets and supporting functionals]
Let $p$ be sublinear on a real vector space.
\begin{enumerate}
  \item The set $\{x:p(x)\le1\}$ is convex.
  \item For every $x_0$ there is a linear functional $f$ such that
        $f\le p$ on $X$ and $f(x_0)=p(x_0)$.
\end{enumerate}
\end{noteexercise}

\begin{proof}
Convexity is immediate from positive homogeneity and subadditivity.  Define
$f_0(\alpha x_0)=\alpha p(x_0)$.  For $\alpha<0$, the inequality
$p(x_0)+p(-x_0)\ge0$ gives
$f_0(\alpha x_0)\le p(\alpha x_0)$.  Hahn--Banach extends $f_0$.
\end{proof}
\sourcepage{70}

\begin{noteexercise}[Hyperplanes]
A subspace $H$ of a vector space $X$ has codimension one if and only if
$H=\Ker f$ for some nonzero linear functional $f$.  If two nonzero
functionals have the same kernel, they are scalar multiples of each other.
\end{noteexercise}

\begin{proof}
If $X/H$ is one-dimensional, compose the quotient map with a nonzero linear
functional on $X/H$.  Conversely, the first isomorphism theorem gives
$X/\Ker f\cong\Ran f$, which is one-dimensional.  The proportionality claim
follows from the unique decomposition relative to a vector outside the common
kernel.
\end{proof}
\sourcepage{71}

\section{The dual of $C([a,b])$}

\begin{notetheorem}[Riesz--Stieltjes representation]
Every bounded linear functional $F$ on $C([a,b])$ has a representation
\[
  F(f)=\int_a^b f(t)\,\mathrm{d}\alpha(t),
\]
where $\alpha$ is a function of bounded variation on $[a,b]$.  With the
normalization $\alpha(a)=0$ and right continuity on $(a,b)$, the representing
function is unique and
\[
  \|F\|=\Var_{[a,b]}(\alpha).
\]
Equivalently, $C([a,b])^*$ is isometrically the space of finite signed or
complex regular Borel measures on $[a,b]$.
\end{notetheorem}

\begin{noteremark}
The notation $C([a,b])^*\cong BV([a,b])$ requires a normalization: adding a
constant to $\alpha$ does not change the Riemann--Stieltjes integral.  The
measure formulation avoids this ambiguity.
\end{noteremark}
\sourcepage{72}

\section{Adjoint operators and annihilators}

\begin{notedefinition}[Adjoint]
Let $T\in\Bop{X}{Y}$.  The adjoint is the bounded linear map
\[
  T^*\colon Y^*\to X^*,
  \qquad
  (T^*g)(x)=g(Tx).
\]
\end{notedefinition}

\begin{notetheorem}[Norm of the adjoint]
For every bounded operator $T$,
\[
  \|T^*\|=\|T\|.
\]
\end{notetheorem}

\begin{proof}
For $g\in Y^*$,
$\|T^*g\|\le\|g\|\|T\|$, so $\|T^*\|\le\|T\|$.  Conversely, given
$\varepsilon>0$, choose $x\in S_X$ with
$\|Tx\|>\|T\|-\varepsilon$.  Hahn--Banach gives $g\in S_{Y^*}$ with
$g(Tx)=\|Tx\|$.  Hence
$\|T^*\|\ge\|T^*g\|\ge|T^*g(x)|>\|T\|-\varepsilon$.
\end{proof}
\sourcepage{73}

\begin{notedefinition}[Annihilators]
For $M\subseteq X$, define
\[
  M^\perp=\{f\in X^*:f(x)=0\text{ for every }x\in M\}.
\]
For $N\subseteq X^*$, define
\[
  {}^\perp N=\{x\in X:f(x)=0\text{ for every }f\in N\}.
\]
\end{notedefinition}

\begin{notetheorem}[Range--kernel annihilator identities]
For $T\in\Bop{X}{Y}$,
\[
  \overline{\Ran T}={}^\perp\Ker T^*,
  \qquad
  \overline{\Ran T^*}^{\;w^*}=(\Ker T)^\perp.
\]
In particular,
\[
  (\Ran T)^\perp=\Ker T^*.
\]
\end{notetheorem}

\begin{proof}
A functional $g\in Y^*$ annihilates $\Ran T$ exactly when
$g(Tx)=0$ for every $x$, that is, exactly when $T^*g=0$.  Taking
preannihilators and using the Hahn--Banach closure theorem gives the first
identity.  The second is the dual statement with weak-$*$ closure.
\end{proof}
\sourcepage{74}

\begin{notedefinition}[Canonical embedding]
For a normed space $X$, define
\[
  J\colon X\to X^{**},
  \qquad
  (Jx)(f)=f(x).
\]
Then $J$ is a linear isometry.
\end{notedefinition}

\begin{proof}
The estimate $\|Jx\|\le\|x\|$ is immediate.  A norming functional
$f\in S_{X^*}$ with $f(x)=\|x\|$ gives the reverse inequality.
\end{proof}
\sourcepage{75}

\begin{notedefinition}[Reflexive space]
A normed space $X$ is reflexive if the canonical embedding is onto:
$J(X)=X^{**}$.
\end{notedefinition}

\begin{notetheorem}
Every reflexive normed space is Banach, and every finite-dimensional normed
space is reflexive.
\end{notetheorem}

\begin{proof}
The bidual is Banach, and $J(X)$ is closed precisely when $X$ is complete.
If $J$ is onto, $X$ is isometric to the Banach space $X^{**}$.  In finite
dimension, $\dim X=\dim X^{**}$ and the injective map $J$ is therefore onto.
\end{proof}
\sourcepage{76}

\begin{noteexercise}[Annihilators separate closed subspaces]
Distinct closed subspaces $Y_1,Y_2$ of a normed space have distinct
annihilators.  A subset $M\subseteq X$ is total, meaning
$\overline{\Span M}=X$, if and only if every $f\in X^*$ vanishing on $M$ is
zero.
\end{noteexercise}

\begin{proof}
If $x\in Y_1\setminus Y_2$, separate $x$ from the closed subspace $Y_2$.
The resulting functional belongs to $Y_2^\perp$ but not $Y_1^\perp$.  The
totality criterion is the annihilator criterion for density.
\end{proof}
\sourcepage{77}

\section{Baire category and its consequences}

\begin{notetheorem}[Baire category theorem]
Let $X$ be a complete metric space.
\begin{enumerate}
  \item $X$ is not a countable union of closed sets with empty interior.
  \item The intersection of countably many dense open subsets of $X$ is
        dense, and in particular nonempty.
\end{enumerate}
\end{notetheorem}

\begin{proof}
It is enough to prove the second statement.  Let $(G_n)$ be dense and open and
let $U$ be a nonempty open set.  Choose a closed ball
$\overline B(x_1,r_1)\subseteq U\cap G_1$.  Inductively choose
\[
  \overline B(x_{n+1},r_{n+1})
  \subseteq B(x_n,r_n)\cap G_{n+1},
  \qquad r_{n+1}<\frac12r_n.
\]
The centres form a Cauchy sequence and converge to some $x\in X$.  The nesting
puts $x$ in every $G_n$ and in $U$.  Taking complements gives the first
statement.
\end{proof}
\sourcepage{78}

\begin{notetheorem}[Cantor intersection theorem]
Let $(F_n)$ be nonempty closed subsets of a complete metric space such that
\[
  F_1\supseteq F_2\supseteq\cdots,
  \qquad
  \diam(F_n)\longrightarrow0.
\]
Then $\bigcap_nF_n$ consists of exactly one point.
\end{notetheorem}

\begin{proof}
Choose $x_n\in F_n$.  If $m,n\ge N$, then
$d(x_m,x_n)\le\diam(F_N)$, so $(x_n)$ is Cauchy.  Let $x_n\to x$.  For every
$N$, all sufficiently late terms lie in $F_N$, and closedness gives
$x\in F_N$.  Two points in the intersection have distance at most
$\diam(F_N)$ for every $N$, hence are equal.
\end{proof}

\begin{notetheorem}[A closed-set form of Baire]
If $X$ is complete and $X=\bigcup_{n=1}^{\infty}F_n$ with each $F_n$ closed,
then some $F_n$ has nonempty interior.  More generally, if an open set
$U\subseteq\bigcup_nF_n$, then some $F_n$ has nonempty interior relative to
$U$.
\end{notetheorem}
\sourcepages{79--80}

\begin{notetheorem}[No countable Hamel basis]
An infinite-dimensional Banach space has no countable Hamel basis.
Consequently, the vector space of polynomials cannot be complete under any
norm for which its usual countable monomial basis is a Hamel basis.
\end{notetheorem}

\begin{proof}
If $(e_n)$ were a Hamel basis, then
\[
  X=\bigcup_{n=1}^{\infty}\Span\{e_1,\ldots,e_n\}.
\]
Each finite-dimensional span is closed and has empty interior: a subspace
with nonempty interior is the whole space.  This contradicts Baire's theorem.
\end{proof}
\sourcepage{81}

\begin{notetheorem}[Cantor's theorem characterizes completeness]
A metric space $X$ is complete if and only if every nested sequence of
nonempty closed sets with diameters tending to zero has nonempty intersection.
\end{notetheorem}

\begin{proof}
Only the converse needs proof.  For a Cauchy sequence $(x_n)$ set
\[
  F_n=\overline{\{x_k:k\ge n\}}.
\]
The sets are nonempty, closed, nested, and their diameters tend to zero.
A point in their intersection is the limit of $(x_n)$.
\end{proof}
\sourcepage{82}

\begin{notetheorem}[Existence of nowhere differentiable functions]
There are continuous real-valued functions on $[0,1]$ that are differentiable
at no point.
\end{notetheorem}

\begin{proof}
Work in the Banach space $C([0,1])$.  For $m,n\in\mathbb N$, let $F_{m,n}$ be
the set of functions $f$ for which there is $x\in[0,1]$ satisfying
\[
  |f(y)-f(x)|\le m|y-x|
  \quad\text{whenever }0<|y-x|<1/n.
\]
Each $F_{m,n}$ is closed: for uniformly convergent $f_k\to f$, choose witness
points $x_k$, pass to a convergent subsequence, and pass the inequality to the
limit.  Each $F_{m,n}$ has empty interior.  Indeed, approximate the centre of
any ball uniformly by a polygonal function and add a sufficiently small,
high-frequency sawtooth whose slopes have absolute value larger than
$m$ plus the finitely many background slopes.  The perturbed function remains
in the ball but has no point satisfying the displayed local Lipschitz bound.

If $f$ is differentiable at $x$, then for some $m,n$ it belongs to $F_{m,n}$.
Thus the set of functions differentiable somewhere is contained in the
countable union $\bigcup_{m,n}F_{m,n}$, a meagre set.  Baire's theorem leaves
a continuous function outside this union.
\end{proof}
\sourcepages{83--84}

\begin{notetheorem}[Uniform boundedness principle]
Let $X$ be Banach, let $Y$ be normed, and let
$\mathcal A\subseteq\Bop{X}{Y}$.  If
\[
  \sup_{T\in\mathcal A}\|Tx\|<\infty
  \qquad(x\in X),
\]
then
\[
  \sup_{T\in\mathcal A}\|T\|<\infty.
\]
\end{notetheorem}

\begin{proof}
Set
\[
  F_n=\{x:\sup_{T\in\mathcal A}\|Tx\|\le n\}.
\]
The sets $F_n$ are closed and cover $X$.  Baire gives
$B(x_0,r)\subseteq F_N$ for some $x_0,r,N$.  If $\|h\|<r$, then
$x_0\pm h\in F_N$, so
\[
  2\|Th\|=\|T(x_0+h)-T(x_0-h)\|\le2N.
\]
Taking $h=(r/2)x$ for $\|x\|\le1$ gives
$\|Tx\|\le2N/r$ uniformly in $T$.
\end{proof}
\sourcepage{85}

\begin{notetheorem}[Banach--Steinhaus pointwise-limit theorem]
Let $X$ be Banach, let $Y$ be normed, and let $A_n\in\Bop{X}{Y}$.  Suppose
$A_nx$ converges for every $x\in X$, and define
$Ax=\lim_nA_nx$.  Then $A\in\Bop{X}{Y}$ and
\[
  \|A\|\le\liminf_{n\to\infty}\|A_n\|.
\]
\end{notetheorem}

\begin{proof}
For every $x$, the sequence $(A_nx)$ is bounded, so uniform boundedness gives
$M=\sup_n\|A_n\|<\infty$.  Linearity passes to the limit and
$\|Ax\|\le M\|x\|$.  For $\|x\|=1$,
\[
  \|Ax\|=\lim_n\|A_nx\|
  \le\liminf_n\|A_n\|;
\]
take the supremum over the unit sphere.
\end{proof}
\sourcepage{86}

\begin{noteexample}[The polynomial space is not complete]
Let $P([0,1])$ be the polynomial space with the supremum norm.  The Taylor
polynomials
\[
  p_n(x)=\sum_{k=0}^{n}\frac{x^k}{k!}
\]
converge uniformly on $[0,1]$ to $e^x$.  Thus $(p_n)$ is Cauchy in the
supremum norm, but its limit is not a polynomial.  Hence $P([0,1])$ is not
complete.  This direct argument also avoids the source's attempted use of
coefficient or derivative functionals, which are not bounded on the full
polynomial space with the supremum norm.
\end{noteexample}
\sourcepages{86--87}

\begin{noteexercise}[Two applications of uniform boundedness]
\begin{enumerate}
  \item If $x_n\in X$ and $f(x_n)$ is bounded for every $f\in X^*$, then
        $(x_n)$ is norm bounded.
  \item For Banach spaces $X,Y$ and $T_n\in\Bop{X}{Y}$, the following are
        equivalent:
        \begin{enumerate}
          \item $(\|T_n\|)$ is bounded;
          \item $(\|T_nx\|)$ is bounded for every $x\in X$;
          \item $(g(T_nx))$ is bounded for every $x\in X$ and $g\in Y^*$.
        \end{enumerate}
\end{enumerate}
\end{noteexercise}

\begin{proof}
For (1), apply uniform boundedness to the canonical evaluations
$Jx_n\in X^{**}$ and use $\|Jx_n\|=\|x_n\|$.  For (2), only the last
implication is nontrivial.  First apply uniform boundedness on $Y^*$ to the
functionals $g\mapsto g(T_nx)$ for each fixed $x$, obtaining pointwise norm
boundedness in $Y$; then apply it on $X$ to the operators $T_n$.
\end{proof}
\sourcepage{88}

\subsection{Quadrature, interpolation, and Fourier sums}

\begin{notetheorem}[P\'olya's theorem for quadrature formulas]
Let $w\in L^1(0,1)$ and let
\[
  L_n(f)=\sum_{j=0}^{m_n}w_j^{(n)}f(x_j^{(n)}),
  \qquad f\in C([0,1]),
\]
where, for each $n$, the nodes $x_j^{(n)}\in[0,1]$ are distinct.  Assume
\[
  L_n(p)\longrightarrow\int_0^1p(t)w(t)\dd t
  \qquad\text{for every polynomial }p.
\]
Then the same convergence holds for every $f\in C([0,1])$ if and only if
\[
  \sup_n\sum_j|w_j^{(n)}|<\infty.
\]
\end{notetheorem}

\begin{proof}
The functional norm is
$\|L_n\|=\sum_j|w_j^{(n)}|$: the upper bound is immediate, and a continuous
piecewise-linear function can be chosen with prescribed signs at the finitely
many nodes.  If the norms are uniformly bounded, approximate $f$ uniformly by
a polynomial $p$ and estimate
\[
 |L_n(f)-\textstyle\int fw|
 \le\|L_n\|\|f-p\|_\infty
     +|L_n(p)-\textstyle\int pw|
     +\|w\|_1\|p-f\|_\infty.
\]
Conversely, convergence for every $f$ makes $(L_n(f))$ pointwise bounded, so
uniform boundedness yields $\sup_n\|L_n\|<\infty$.
\end{proof}
\sourcepages{89--90}

\begin{notetheorem}[Norm of the Lagrange interpolation operator]
For distinct nodes $0\le t_0<\cdots<t_n\le1$, let
\[
  (L_nf)(t)=\sum_{k=0}^nf(t_k)p_k(t),
  \qquad
  p_k(t)=\prod_{j\ne k}\frac{t-t_j}{t_k-t_j}.
\]
Then
\[
  \|L_n\|
  =\sup_{t\in[0,1]}\sum_{k=0}^n|p_k(t)|.
\]
\end{notetheorem}

\begin{proof}
The triangle inequality gives the upper bound.  Choose $t^*$ attaining the
maximum and a continuous piecewise-linear function $f$ with
$\|f\|_\infty=1$ and
$f(t_k)=\sign p_k(t^*)$ (using arbitrary phases in the complex case).  Then
$|(L_nf)(t^*)|=\sum_k|p_k(t^*)|$.
\end{proof}
\sourcepages{91--92}

\begin{notetheorem}[Fourier partial sums and the Dirichlet kernel]
On the Banach space of continuous $2\pi$-periodic functions, the $n$th Fourier
partial-sum operator satisfies
\[
  (S_nf)(t)=\frac1{2\pi}\int_{-\pi}^{\pi}f(t-u)D_n(u)\dd u,
  \qquad
  D_n(u)=\frac{\sin((n+\tfrac12)u)}{\sin(u/2)},
\]
and
\[
  \|S_n\|=\frac1{2\pi}\int_{-\pi}^{\pi}|D_n(u)|\dd u.
\]
Moreover, $\|S_n\|\to\infty$ at logarithmic order.  Hence there exists a
continuous periodic function whose Fourier partial sums diverge at a point.
\end{notetheorem}

\begin{proof}
The convolution formula gives the upper bound.  For the reverse bound,
approximate the phase function $u\mapsto\sign(D_n(-u))$ by continuous
periodic functions, bounded in modulus by one, in the weighted space
$L^1(|D_n(u)|\dd u)$.  On each interval where the numerator has fixed sign
and $|u|$ is small, $|\sin(u/2)|\le |u|/2$.  Summing the resulting harmonic
lower bounds yields $\|S_n\|\to\infty$.  If every continuous $f$ had
bounded partial sums at zero, uniform boundedness would force
$\sup_n\|S_n\|<\infty$, a contradiction.
\end{proof}
\sourcepages{93--94}

\begin{notetheorem}[Separate continuity of bilinear maps]
Let $X$ and $Y$ be Banach spaces, let $Z$ be normed, and let
$B\colon X\times Y\to Z$ be bilinear.  If $B(\,\cdot\,,y)$ is bounded for
every fixed $y$ and $B(x,\,\cdot\,)$ is bounded for every fixed $x$, then
there is $C<\infty$ such that
\[
  \|B(x,y)\|\le C\|x\|\|y\|.
\]
Thus $B$ is jointly continuous.
\end{notetheorem}

\begin{proof}
For $y\in Y$, define $A_yx=B(x,y)$.  For each fixed $x$, the family
$\{A_y:\|y\|\le1\}$ is pointwise bounded because
$y\mapsto B(x,y)$ is bounded.  Uniform boundedness gives
$\sup_{\|y\|\le1}\|A_y\|=C<\infty$, which is the desired estimate.
\end{proof}

\begin{noteexample}[Completeness is essential]
Let $c_{00}$ carry the supremum norm and define
\[
  B(x,y)=\sum_{k=1}^{\infty}x_ky_k.
\]
The sum is finite.  For fixed $x$ or fixed $y$, the resulting functional is
bounded, but $B$ is not jointly bounded: taking $x=y=(1,\ldots,1,0,\ldots)$
with $n$ ones gives $B(x,y)=n$ while both sup norms equal one.
\end{noteexample}
\sourcepages{95--96}

\begin{notecorollary}[Positive quadrature weights]
In P\'olya's theorem, if every $w_j^{(n)}\ge0$, convergence on polynomials
already implies convergence on all continuous functions.
\end{notecorollary}

\begin{proof}
Applying polynomial convergence to the constant polynomial $1$ gives
\[
  \sum_j|w_j^{(n)}|=\sum_jw_j^{(n)}=L_n(1)
  \longrightarrow\int_0^1w(t)\dd t.
\]
The norms are therefore uniformly bounded, and P\'olya's theorem applies.
\end{proof}
\sourcepage{97}

\section{Weak and weak-star convergence}

\begin{notedefinition}[Strong and weak convergence]
A sequence $(x_n)$ in a normed space $X$ converges strongly to $x$ if
$\|x_n-x\|\to0$.  It converges weakly to $x$, written
$x_n\weakto x$, if
\[
  f(x_n)\longrightarrow f(x)
  \qquad(f\in X^*).
\]
\end{notedefinition}

\begin{notelemma}[Basic facts]
If $x_n\weakto x$, then:
\begin{enumerate}
  \item the weak limit is unique;
  \item every subsequence converges weakly to $x$;
  \item $(x_n)$ is norm bounded;
  \item $\|x\|\le\liminf_n\|x_n\|$.
\end{enumerate}
\end{notelemma}

\begin{proof}
Functionals separate points, giving uniqueness.  The subsequence statement is
immediate.  For boundedness, regard $Jx_n$ as functionals on the Banach space
$X^*$.  They are pointwise bounded, so uniform boundedness gives
$\sup_n\|Jx_n\|=\sup_n\|x_n\|<\infty$.  A norming functional for $x$ gives
\[
  \|x\|=|f(x)|=\lim_n|f(x_n)|
  \le\liminf_n\|x_n\|.
\]
\end{proof}
\sourcepages{98--99}

\begin{notetheorem}[Strong versus weak convergence]
Strong convergence implies weak convergence.  The converse holds in
finite-dimensional normed spaces, but fails in every infinite-dimensional
Hilbert space: an orthonormal sequence $(e_n)$ satisfies $e_n\weakto0$ while
$\|e_n\|=1$.
\end{notetheorem}

\begin{proof}
For strong convergence,
$|f(x_n)-f(x)|\le\|f\|\|x_n-x\|$.  In finite dimension, weak convergence of
the coordinate functionals gives convergence of all coordinates and hence
norm convergence.  In a Hilbert space, Riesz representation writes every
$f$ as $f(x)=\langle x,z\rangle$; Bessel's inequality gives
$\langle e_n,z\rangle\to0$.
\end{proof}

\begin{noteremark}[Schur's theorem]
In $\ell^1$, weak convergence of sequences implies norm convergence.  This is
a special property of $\ell^1$, not a general feature of infinite-dimensional
Banach spaces.
\end{noteremark}
\sourcepages{99--100}

\begin{notelemma}[Testing weak convergence on a total set]
Let $(x_n)$ be bounded in $X$ and let $M\subseteq X^*$ have norm-dense linear
span in $X^*$.  If $f(x_n)\to f(x)$ for every $f\in M$, then
$x_n\weakto x$.
\end{notelemma}

\begin{proof}
The convergence extends by linearity to $\Span M$.  If
$\|x_n\|,\|x\|\le C$ and $g\in X^*$, choose $f\in\Span M$ with
$\|g-f\|<\varepsilon/(3C)$.  Then
\[
 |g(x_n-x)|
 \le |(g-f)(x_n)|+|f(x_n-x)|+|(f-g)(x)|<\varepsilon
\]
for all sufficiently large $n$.
\end{proof}
\sourcepage{100}

\begin{notetheorem}[Weak convergence in $\ell^p$]
Let $1<p<\infty$.  A sequence $x^{(n)}=(x_k^{(n)})$ in $\ell^p$ converges
weakly to $x=(x_k)$ if and only if
\begin{enumerate}
  \item $(\|x^{(n)}\|_p)$ is bounded;
  \item $x_k^{(n)}\to x_k$ for every coordinate $k$.
\end{enumerate}
\end{notetheorem}

\begin{proof}
Weak convergence gives both conditions.  Conversely, coordinate evaluations
span the finitely supported sequences, which are dense in
$(\ell^p)^*=\ell^q$.  Apply the total-set criterion.
\end{proof}

\begin{notetheorem}[Radon--Riesz property in uniformly convex spaces]
If $X$ is uniformly convex, $x_n\weakto x$, and
$\|x_n\|\to\|x\|$, then $x_n\to x$ in norm.
\end{notetheorem}

\begin{proof}
The case $x=0$ is immediate.  Otherwise put
$u_n=x_n/\|x_n\|$ and $u=x/\|x\|$.  Then $u_n\weakto u$ and
\[
  \left\|\frac{u_n+u}{2}\right\|
  \ge \operatorname{Re}f\!\left(\frac{u_n+u}{2}\right)\longrightarrow1
\]
for a norming functional $f$ of $u$.  Uniform convexity forces
$\|u_n-u\|\to0$, and the convergence of norms gives $\|x_n-x\|\to0$.
\end{proof}
\sourcepage{101}

\begin{notetheorem}[Weak convergence and operators]
Let $A\in\Bop{X}{Y}$.
\begin{enumerate}
  \item If $x_n\weakto x$, then $Ax_n\weakto Ax$.
  \item If $A$ is compact and $x_n\weakto x$, then
        $\|Ax_n-Ax\|\to0$.
  \item If $B\colon X\times Y\to Z$ is bounded bilinear,
        $x_n\to x$ strongly, and $y_n\weakto y$, then
        $B(x_n,y_n)\weakto B(x,y)$; when $Z=\mathbb K$, the convergence is
        scalar.
\end{enumerate}
\end{notetheorem}

\begin{proof}
For (1), $g(Ax_n)=(A^*g)(x_n)$.  For (2), weak convergence makes $(x_n)$
bounded.  Every subsequence of $(Ax_n)$ has a norm-convergent further
subsequence, and (1) forces its limit to be $Ax$; hence the whole sequence
converges in norm.  For (3), $(y_n)$ is bounded and
\[
  \|B(x_n,y_n)-B(x,y_n)\|
  \le\|B\|\|x_n-x\|\|y_n\|\to0,
\]
while $B(x,\cdot)$ is a bounded linear map and therefore preserves weak
convergence.
\end{proof}
\sourcepage{102}

\begin{notedefinition}[Weak topology and weak-$*$ topology]
The weak topology on $X$ is the coarsest topology making every member of
$X^*$ continuous.  The weak-$*$ topology on $X^*$ is the coarsest topology
making every evaluation $f\mapsto f(x)$, $x\in X$, continuous.
\end{notedefinition}

\begin{notetheorem}[Finite dimensions and sequences]
In finite dimension, the weak and norm topologies coincide.  In an
infinite-dimensional normed space the weak topology is strictly coarser.
For sequences,
\[
  x_n\weakto x
  \quad\Longleftrightarrow\quad
  f(x_n)\to f(x)\text{ for all }f\in X^*.
\]
\end{notetheorem}
\sourcepage{103}

\begin{notedefinition}[Weak-$*$ convergence of functionals]
A sequence $(f_n)$ in $X^*$ converges weak-$*$ to $f$ if
\[
  f_n(x)\longrightarrow f(x)
  \qquad(x\in X).
\]
Weak convergence in $X^*$ implies weak-$*$ convergence.  If $X$ is reflexive,
the two notions agree on $X^*$.
\end{notedefinition}

\begin{notetheorem}[Strong operator limits]
Let $X$ be Banach, let $Y$ be normed, and let
$T_n\in\Bop{X}{Y}$.  If $T_nx$ converges in $Y$ for every $x\in X$, then the
pointwise limit $Tx=\lim_nT_nx$ belongs to $\Bop{X}{Y}$ and
$\sup_n\|T_n\|<\infty$.
\end{notetheorem}

\begin{proof}
The family $(T_n)$ is pointwise bounded, so uniform boundedness gives a common
operator-norm bound.  Linearity and boundedness pass to the pointwise limit.
\end{proof}
\sourcepage{104}

\begin{notedefinition}[Operator convergence]
For $T_n,T\in\Bop{X}{Y}$:
\begin{itemize}
  \item $T_n\to T$ uniformly if $\|T_n-T\|\to0$;
  \item $T_n\to T$ strongly if $T_nx\to Tx$ for every $x\in X$;
  \item $T_n\to T$ weakly if $T_nx\weakto Tx$ for every $x\in X$.
\end{itemize}
\end{notedefinition}

\begin{notetheorem}[Criterion for strong operator convergence]
Let $X,Y$ be Banach spaces and let $(T_n)\subseteq\Bop{X}{Y}$.  The sequence
converges strongly if and only if
\begin{enumerate}
  \item $\sup_n\|T_n\|<\infty$;
  \item for every $x$ in a total subset $M\subseteq X$, the sequence
        $(T_nx)$ is Cauchy in $Y$.
\end{enumerate}
\end{notetheorem}

\begin{proof}
Only sufficiency needs proof.  Cauchy convergence extends by linearity to
$\Span M$.  Given $x\in X$, approximate it by $y\in\Span M$ and use the
uniform bound to control
$\|(T_n-T_m)(x-y)\|$, while $(T_ny)$ is Cauchy.  Completeness of $Y$ supplies
the pointwise limit.
\end{proof}

\begin{notecorollary}
A sequence $(f_n)\subseteq X^*$ converges weak-$*$ if and only if it is norm
bounded and $f_n(x)$ is Cauchy for every $x$ in a total subset of $X$.
\end{notecorollary}
\sourcepage{105}

\begin{notetheorem}[Reflexivity results recorded in the source]
The following standard results are used later.
\begin{enumerate}
  \item Every finite-dimensional normed space is reflexive.
  \item Every Hilbert space is reflexive.
  \item Closed subspaces and finite products of reflexive Banach spaces are
        reflexive.
  \item $L^p(\Omega)$ is reflexive for $1<p<\infty$.
  \item Every uniformly convex Banach space is reflexive
        (Milman--Pettis theorem).
  \item A Banach space is reflexive if and only if every bounded sequence has
        a weakly convergent subsequence (Eberlein--\v{S}mulian plus weak
        compactness of the unit ball).
\end{enumerate}
\end{notetheorem}

\begin{proof}[Editorial note on depth]
Milman--Pettis and Eberlein--\v{S}mulian are deep theorems.  The source records
them without complete proofs; they are therefore stated accurately here and
not represented as consequences of the short handwritten sketches.
\end{proof}
\sourcepage{106}

\begin{notetheorem}[Completeness and the canonical image]
For a normed space $X$, the canonical image $J(X)$ is closed in $X^{**}$ if
and only if $X$ is Banach.
\end{notetheorem}

\begin{proof}
If $X$ is Banach, its isometric image is complete and hence closed.  Conversely,
if $J(X)$ is closed, it is a Banach subspace of $X^{**}$, and the isometry
$J\colon X\to J(X)$ makes $X$ complete.
\end{proof}
\sourcepage{107}

\begin{notelemma}[Dense testing in the dual]
Let $(x_n)$ be bounded in a normed space $X$ and let $Y\subseteq X^*$ be
norm dense.  If $y(x_n)\to y(x)$ for every $y\in Y$, then $x_n\weakto x$.
\end{notelemma}

\begin{notetheorem}[Bounded almost-everywhere convergence in $L^p$]
Let $\Omega\subseteq\mathbb R^d$ be open, let $1<p<\infty$, and let
$(f_n)$ be bounded in $L^p(\Omega)$.  If $f_n\to f$ almost everywhere and
$f\in L^p(\Omega)$, then $f_n\weakto f$ in $L^p(\Omega)$.
\end{notetheorem}

\begin{proof}
Let $q$ be conjugate to $p$.  Given $g\in L^q(\Omega)$, first approximate it
in $L^q$ by a bounded function $g_0$ supported on a set of finite measure.
The common $L^p$ bound controls the error uniformly in $n$.  On the support of
$g_0$, almost-everywhere convergence together with the common $L^p$ bound
implies $f_n-f\to0$ in $L^1$: use Egorov's theorem off a set of arbitrarily
small measure and H\"older's inequality on that exceptional set.  Therefore
$\int g_0(f_n-f)\to0$, and the approximation gives
$\int g(f_n-f)\to0$.
\end{proof}

\begin{noteexample}[The alternating-cell sequence in $L^2(0,1)$]
For $n\ge1$ and $0\le j\le n-1$, define
\[
 f_n(t)=
 \begin{cases}
   0,&\displaystyle \frac jn\le t<\frac jn+\frac1{2n},\\[4pt]
   1,&\displaystyle \frac jn+\frac1{2n}\le t<\frac{j+1}{n}.
 \end{cases}
\]
Then $f_n\weakto\tfrac12$ in $L^2(0,1)$.  Indeed, the integrals of
$f_n-\tfrac12$ over every interval $[0,t]$ tend to zero, and finite linear
combinations of interval indicators are dense in $L^2(0,1)$; the sequence is
bounded, so the dense-testing lemma applies.  The convergence is not strong,
because
\[
  \left\|f_n-\frac12\right\|_2=\frac12
  \qquad(n\ge1).
\]
\end{noteexample}
\sourcepages{108--109}

\begin{notetheorem}[Weak lower semicontinuity of the norm]
If $x_n\weakto x$, then
\[
  \|x\|\le\liminf_{n\to\infty}\|x_n\|.
\]
\end{notetheorem}

\begin{proof}
Choose $f\in X^*$ with $\|f\|=1$ and $f(x)=\|x\|$.  Then
$|f(x_n)|\le\|x_n\|$ and $f(x_n)\to f(x)$, which gives the inequality.
\end{proof}

\begin{noteexercise}[Weak limits lie in closed spans]
If $x_n\in Y$ for a linear subspace $Y\subseteq X$ and $x_n\weakto x$, then
$x\in\overline Y$.
\end{noteexercise}

\begin{proof}
If $x\notin\overline Y$, Hahn--Banach gives $f\in X^*$ vanishing on $Y$ but
not at $x$, contradicting $f(x_n)\to f(x)$.
\end{proof}
\sourcepage{109}

\begin{notetheorem}[Weak-$*$ sequential compactness in the separable case]
Let $X$ be separable and let $(f_n)$ be a bounded sequence in $X^*$.  Then
$(f_n)$ has a weak-$*$ convergent subsequence.
\end{notetheorem}

\begin{proof}
Choose a dense sequence $(x_j)$ in $X$.  Repeatedly extract subsequences so
that $f_n(x_j)$ converges for each $j$, then take the diagonal subsequence.
The pointwise limit on the dense span is bounded by the common norm bound and
therefore extends uniquely to a member of $X^*$.  Approximation from the dense
span shows weak-$*$ convergence on all of $X$.
\end{proof}

\begin{noteremark}
This is the sequential form of Banach--Alaoglu for separable preduals.  In
nonseparable settings, weak-$*$ compactness need not be sequential compactness.
\end{noteremark}
\sourcepage{110}

\begin{notetheorem}[Toeplitz limit theorem]
Let $A=(a_{nk})_{n,k\ge1}$ be an infinite matrix and define
\[
  (Ax)_n=\sum_{k=1}^{\infty}a_{nk}x_k
\]
whenever the sum is meaningful.  The method is regular, meaning it preserves
limits of convergent scalar sequences, if and only if
\[
  \sup_n\sum_{k=1}^{\infty}|a_{nk}|<\infty,
  \qquad
  a_{nk}\longrightarrow0\quad(n\to\infty)\text{ for each fixed }k,
\]
and
\[
  \sum_{k=1}^{\infty}a_{nk}\longrightarrow1.
\]
\end{notetheorem}

\begin{proof}[Proof sketch]
The necessity follows by testing the method on the constant sequence and on
the coordinate sequences, while uniform boundedness on the Banach space $c$
of convergent sequences gives the uniform row-$\ell^1$ bound.  Conversely,
write $x_k=L+(x_k-L)$, split the remainder into a finite head and a uniformly
small tail, and use the three conditions.
\end{proof}
\sourcepage{111}

\section{Open mapping and bounded inverse theorems}

\begin{notetheorem}[Open mapping theorem]
Let $X$ and $Y$ be Banach spaces and let $T\in\Bop{X}{Y}$ be surjective.
Then $T$ is an open map.
\end{notetheorem}

\begin{proof}
It is enough to show that $T(B_X)$ contains a neighbourhood of $0$.  Since
\[
  Y=\bigcup_{n=1}^{\infty}T(nB_X),
\]
Baire's theorem gives $n$ such that $\overline{T(nB_X)}$ has nonempty
interior.  Taking differences and rescaling yields a number $\delta>0$ with
\[
  B_Y(0,\delta)\subseteq\overline{T(B_X)}.
\]
After rescaling, assume $B_Y(0,1)\subseteq\overline{T(B_X)}$.  Given
$y\in B_Y(0,1/2)$, choose inductively $x_k\in2^{-k}B_X$ so that
\[
  \left\|y-T\sum_{j=1}^{k}x_j\right\|<2^{-k-1}.
\]
The series $\sum_kx_k$ converges to $x\in B_X$, and continuity gives
$Tx=y$.  Thus a ball around $0$ lies in $T(B_X)$; translation and scaling
prove openness.
\end{proof}
\sourcepages{112--113}

\begin{notecorollary}[Bounded inverse theorem]
A bounded bijective linear map between Banach spaces has a bounded inverse.
\end{notecorollary}

\begin{proof}
The inverse is continuous because the original map is open.
\end{proof}

\begin{noteexample}[A priori estimate for a boundary-value problem]
Let
\[
  L\colon C^2_0([0,1])\to C([0,1]),
  \qquad
  Lu=a u''+b u'+cu,
\]
where $C^2_0([0,1])=\{u\in C^2([0,1]):u(0)=u(1)=0\}$ carries the norm
$\|u\|_\infty+\|u'\|_\infty+\|u''\|_\infty$.  If the two-point
boundary-value problem $Lu=f$ has a unique solution for every
$f\in C([0,1])$, then $L$ is a bounded bijection between Banach spaces.
Therefore there is $C>0$ such that
\[
  \|u\|_\infty+\|u'\|_\infty+\|u''\|_\infty
  \le C\|Lu\|_\infty.
\]
\end{noteexample}
\sourcepage{113}

\begin{notetheorem}[Equivalence of complete norms]
Let $\|\cdot\|_1$ and $\|\cdot\|_2$ be two norms on the same vector space
$X$, and assume both make $X$ Banach.  If
\[
  \|x\|_2\le C\|x\|_1
\]
for all $x$, then the norms are equivalent.
\end{notetheorem}

\begin{proof}
The identity map from $(X,\|\cdot\|_1)$ to $(X,\|\cdot\|_2)$ is a bounded
bijection.  The bounded inverse theorem gives
$\|x\|_1\le C'\|x\|_2$.
\end{proof}
\sourcepage{114}

\begin{notetheorem}[Stability of surjectivity]
Let $X,Y$ be Banach spaces and let $A\in\Bop{X}{Y}$ be surjective.  Then
there is $\varepsilon>0$ such that every
$\widetilde A\in\Bop{X}{Y}$ with
$\|\widetilde A-A\|<\varepsilon$ is surjective.
\end{notetheorem}

\begin{proof}
The open mapping theorem gives $c>0$ such that every $y\in Y$ has a preimage
$x$ with $Ax=y$ and $\|x\|\le c\|y\|$.  Choose
$\varepsilon<1/c$.  For a target $y$, construct $x_0,x_1,\ldots$ with
\[
  Ax_0=y,
  \qquad
  Ax_{n+1}=(A-\widetilde A)x_n,
  \qquad
  \|x_{n+1}\|\le c\|A-\widetilde A\|\|x_n\|.
\]
The series $x=\sum_nx_n$ converges and telescoping gives
$\widetilde Ax=y$.
\end{proof}
\sourcepage{115}

\begin{notetheorem}[Closed range and bounded inverse on the range]
Let $X,Y$ be Banach spaces and let $T\in\Bop{X}{Y}$ be injective.  The
inverse $T^{-1}\colon\Ran T\to X$ is bounded if and only if $\Ran T$ is
closed in $Y$.
\end{notetheorem}

\begin{proof}
If $T^{-1}$ is bounded and $Tx_n\to y$, then $(x_n)$ is Cauchy and converges
to $x$; hence $y=Tx\in\Ran T$.  Conversely, if $\Ran T$ is closed, it is
Banach, and the bounded bijection $T\colon X\to\Ran T$ has bounded inverse.
\end{proof}
\sourcepage{116}

\section{Closed graph theorem and closed operators}

\begin{notetheorem}[Products of Banach spaces]
If $X$ and $Y$ are Banach spaces, then $X\times Y$ is Banach under either of
the equivalent norms
\[
  \|(x,y)\|_1=\|x\|+\|y\|,
  \qquad
  \|(x,y)\|_\infty=\max\{\|x\|,\|y\|\}.
\]
\end{notetheorem}
\sourcepage{117}

\begin{notedefinition}[Closed linear operator]
Let $T\colon\Dom(T)\subseteq X\to Y$ be linear.  Its graph is
\[
  \Graph(T)=\{(x,Tx):x\in\Dom(T)\}\subseteq X\times Y.
\]
The operator is \emph{closed} if its graph is closed.
\end{notedefinition}

\begin{notetheorem}[Closed graph theorem]
Let $X,Y$ be Banach spaces.  An everywhere-defined linear map
$T\colon X\to Y$ is bounded if and only if its graph is closed.
\end{notetheorem}

\begin{proof}
Boundedness plainly implies closedness.  Conversely, if the graph is closed,
it is a Banach space with the graph norm
\[
  \|x\|_T=\|x\|+\|Tx\|.
\]
The projection $S\colon\Graph(T)\to X$, $S(x,Tx)=x$, is a bounded bijection.
The bounded inverse theorem makes $S^{-1}$ bounded, so
\[
  \|x\|+\|Tx\|=\|S^{-1}x\|_T\le C\|x\|.
\]
Thus $T$ is bounded.
\end{proof}
\sourcepages{117--119}

\begin{notetheorem}[Sequential criterion]
A linear operator $T\colon\Dom(T)\subseteq X\to Y$ is closed if and only if
whenever
\[
  x_n\in\Dom(T),
  \qquad
  x_n\to x,
  \qquad
  Tx_n\to y,
\]
one has $x\in\Dom(T)$ and $Tx=y$.
\end{notetheorem}

\begin{notetheorem}[Closedness from a closed domain]
If $\Dom(T)$ is closed in $X$ and $T\colon\Dom(T)\to Y$ is bounded, then
$T$ is closed.  Conversely, if $X$ and $Y$ are Banach and $T$ is closed,
then $\Dom(T)$ is complete for the graph norm, not necessarily for the
original norm.
\end{notetheorem}

\begin{noteexample}[Differentiation]
Let
\[
  D\colon C^1([0,1])\subset C([0,1])\to C([0,1]),
  \qquad Du=u'.
\]
With the supremum norm on both copies of $C([0,1])$, the operator is closed
but unbounded.  Indeed, $u_n\to u$ uniformly and $u_n'\to v$ uniformly imply
\[
  u_n(t)-u_n(0)=\int_0^tu_n'(s)\dd s
  \longrightarrow\int_0^tv(s)\dd s,
\]
so $u\in C^1$ and $u'=v$.  On the other hand,
$u(t)=t^n$ has norm one while $\|u'\|_\infty=n$.
\end{noteexample}
\sourcepage{120}

\begin{noteexercise}[Inverse and compact-set properties]
Let $T\colon\Dom(T)\subseteq X\to Y$ be closed.
\begin{enumerate}
  \item If $T^{-1}$ exists, then $T^{-1}$ is closed.
  \item The image $T(C)$ of a compact set $C\subseteq\Dom(T)$ is closed in
        $Y$.
  \item The inverse image $T^{-1}(K)$ of a compact set $K\subseteq Y$ is
        closed in $X$.
\end{enumerate}
\end{noteexercise}

\begin{proof}
The graph of $T^{-1}$ is obtained from the graph of $T$ by swapping the two
coordinates.  For (2), a convergent sequence $Tx_n\to y$ with $x_n\in C$
has a convergent subsequence $x_{n_k}\to x\in C$; closedness gives $Tx=y$.
For (3), if $x_n\to x$ and $Tx_n\in K$, compactness gives a subsequence
$Tx_{n_k}\to y\in K$, and closedness gives $Tx=y$.
\end{proof}
\sourcepage{121}

\begin{notetheorem}[Compact domain]
Let $K$ be a compact metric space, let $Y$ be a normed space, and let
$T\colon K\to Y$ be a bijection whose graph is closed in $K\times Y$.
Then $T^{-1}\colon T(K)\to K$ is continuous.
\end{notetheorem}

\begin{proof}
If $y_n\to y$ in $T(K)$ and $x_n=T^{-1}y_n$, compactness gives a
convergent subsequence $x_{n_k}\to x\in K$.  Closedness gives $Tx=y$, so
$x=T^{-1}y$.  Every subsequence has a further subsequence converging to the
same point $T^{-1}y$, hence the whole sequence converges.
\end{proof}

\begin{notetheorem}[Closed range from a bounded inverse]
Let $T$ be closed, let $X$ be Banach, and suppose
$T^{-1}\colon\Ran T\to X$ exists and is bounded.  Then $\Ran T$ is closed.
\end{notetheorem}

\begin{proof}
If $y_n=Tx_n\to y$, boundedness of $T^{-1}$ makes $(x_n)$ Cauchy, hence
$x_n\to x\in X$.  Closedness gives $x\in\Dom(T)$ and $Tx=y$.
\end{proof}
\sourcepage{122}

\begin{notetheorem}[Termwise differentiation]
Suppose $u_n\in C^1([0,1])$, the series $\sum_nu_n$ converges uniformly to
$u$, and the derivative series $\sum_nu_n'$ converges uniformly to $v$.
Then $u\in C^1([0,1])$ and $u'=v$.
\end{notetheorem}

\begin{proof}
The partial sums $s_N$ satisfy
$(s_N,Ds_N)\to(u,v)$ in $C([0,1])\times C([0,1])$.  Closedness of the
differentiation operator gives the conclusion.
\end{proof}

\begin{notetheorem}[Closed graph extension criterion]
Let $G\subseteq X\times Y$ be a linear subspace.  There is a linear operator
$T$ with $G=\Graph(T)$ if and only if
\[
  (0,y)\in G\quad\Longrightarrow\quad y=0.
\]
If $G$ is closed, the resulting operator is closed.
\end{notetheorem}

\begin{proof}
The condition ensures that for each $x$ there is at most one $y$ with
$(x,y)\in G$.  Define the domain as the set of such $x$ and set $Tx=y$.
Linearity follows because $G$ is a subspace, and closedness of $G$ is exactly
closedness of the operator.
\end{proof}
\sourcepages{123--124}

\chapter{Hilbert Spaces}

Throughout this chapter, the inner product is linear in its first argument and
conjugate linear in its second.

\section{Inner products and their geometry}

\begin{notedefinition}[Inner product]
An inner product on a real or complex vector space $H$ is a map
$\langle\cdot,\cdot\rangle\colon H\times H\to\mathbb K$ satisfying
\begin{enumerate}
  \item $\langle x,x\rangle\ge0$, with equality only for $x=0$;
  \item $\langle \alpha x+\beta y,z\rangle
        =\alpha\langle x,z\rangle+\beta\langle y,z\rangle$;
  \item $\langle x,y\rangle=\overline{\langle y,x\rangle}$.
\end{enumerate}
The associated norm is
\[
  \|x\|=\sqrt{\langle x,x\rangle}.
\]
A complete inner-product space is a Hilbert space.
\end{notedefinition}

\begin{notetheorem}[Cauchy--Schwarz inequality]
For all $x,y\in H$,
\[
  |\langle x,y\rangle|\le\|x\|\,\|y\|.
\]
Equality holds exactly when $x$ and $y$ are linearly dependent.
\end{notetheorem}

\begin{proof}
If $y\ne0$, positivity of
$\|x-\langle x,y\rangle y/\|y\|^2\|^2$ gives the inequality.  Equality is
equivalent to that orthogonal remainder being zero.
\end{proof}

\begin{notedefinition}[Orthogonality]
Vectors $x,y$ are orthogonal, written $x\perp y$, if
$\langle x,y\rangle=0$.
\end{notedefinition}

\begin{notetheorem}[Pythagorean theorem and orthogonal decomposition]
If $x\perp y$, then
\[
  \|x+y\|^2=\|x\|^2+\|y\|^2.
\]
For $y\ne0$ every $x$ has the decomposition
\[
  x=\frac{\langle x,y\rangle}{\|y\|^2}y
    +\left(x-\frac{\langle x,y\rangle}{\|y\|^2}y\right),
\]
where the second term is orthogonal to $y$.
\end{notetheorem}

\begin{notetheorem}[Triangle inequality]
The inner-product norm satisfies
\[
  \|x+y\|\le\|x\|+\|y\|.
\]
Equality holds exactly when one vector is a nonnegative real multiple of the
other.
\end{notetheorem}

\begin{notetheorem}[Parallelogram identity]
For all $x,y\in H$,
\[
  \|x+y\|^2+\|x-y\|^2=2\|x\|^2+2\|y\|^2.
\]
\end{notetheorem}

\begin{noteexample}[$\ell^p$ is not Hilbert when $p\ne2$]
For $x=(1,1,0,\ldots)$ and $y=(1,-1,0,\ldots)$,
\[
  \|x+y\|_p^2+\|x-y\|_p^2=8,
  \qquad
  2\|x\|_p^2+2\|y\|_p^2=4\cdot2^{2/p}.
\]
The parallelogram identity holds only for $p=2$, so the $\ell^p$ norm is not
induced by an inner product when $p\ne2$.
\end{noteexample}

\begin{notetheorem}[Polarization identities]
A norm induced by an inner product determines that inner product.  Over
$\mathbb R$,
\[
  \langle x,y\rangle
  =\frac14\bigl(\|x+y\|^2-\|x-y\|^2\bigr).
\]
Over $\mathbb C$ with linearity in the first variable,
\[
  \langle x,y\rangle
  =\frac14\sum_{k=0}^{3}i^k\|x+i^ky\|^2.
\]
\end{notetheorem}

\begin{notelemma}[Continuity of the inner product]
If $x_n\to x$ and $y_n\to y$, then
$\langle x_n,y_n\rangle\to\langle x,y\rangle$.
\end{notelemma}

\begin{proof}
By Cauchy--Schwarz,
\[
 |\langle x_n,y_n\rangle-\langle x,y\rangle|
 \le\|x_n-x\|\|y_n\|+\|x\|\|y_n-y\|.
\]
The sequence $(y_n)$ is bounded.
\end{proof}

\begin{notetheorem}[Operator norm through the inner product]
Let $A\in\Bop{H_1}{H_2}$ between Hilbert spaces.  Then
\[
  \|A\|
  =\sup_{x\ne0,\,y\ne0}
    \frac{|\langle Ax,y\rangle|}{\|x\|\|y\|}.
\]
\end{notetheorem}

\begin{proof}
Cauchy--Schwarz gives the upper bound.  For $x$ near norm attaining and
$y=Ax$, the quotient equals $\|Ax\|/\|x\|$.
\end{proof}

\begin{notetheorem}[Isometries fixing the origin]
Let $H_1,H_2$ be real inner-product spaces and let
$A\colon H_1\to H_2$ satisfy $A0=0$ and
$\|Ax-Ay\|=\|x-y\|$.  Then $A$ is linear and preserves inner products.
\end{notetheorem}

\begin{proof}
Taking $y=0$ gives $\|Ax\|=\|x\|$.  The real polarization identity gives
$\langle Ax,Ay\rangle=\langle x,y\rangle$.  Hence
\[
 \|A(x+y)-Ax-Ay\|^2=0,
\]
so $A$ is additive.  Integer and rational homogeneity follow, and continuity
extends homogeneity to all real scalars.
\end{proof}

\begin{noteexample}[Complex caveat]
Complex conjugation on $\mathbb C$ is an isometry fixing $0$ but is not
complex linear.  Thus the preceding statement is a real theorem unless an
additional complex-linearity condition is imposed.
\end{noteexample}

\begin{notetheorem}[Mazur--Ulam theorem]
Every surjective isometry between real normed spaces is affine.  Thus, if
$f(0)=0$, it is real linear.
\end{notetheorem}

\begin{proof}[Proof sketch]
A reflection argument shows that a surjective isometry preserves midpoints:
$f((x+y)/2)=(f(x)+f(y))/2$.  Iterating yields preservation of dyadic convex
combinations, continuity gives all real convex combinations, and translation
by $f(0)$ turns the map into an additive homogeneous map.  The source's
strictly-convex ball-intersection proof covers only a special case; the
statement above is the full theorem.
\end{proof}

\begin{noteexercise}[Vanishing quadratic form]
Let $H$ be a complex inner-product space and let $T\in\Bop{H}{H}$.  If
\[
  \langle Tx,x\rangle=0\qquad(x\in H),
\]
then $T=0$.  This is false over $\mathbb R$.
\end{noteexercise}

\begin{proof}
Apply complex polarization to the sesquilinear form
$b(x,y)=\langle Tx,y\rangle$.  Its diagonal values determine the entire form,
so $b=0$ and $T=0$.  Over $\mathbb R$, a nonzero skew-symmetric rotation of
$\mathbb R^2$ satisfies $\langle Tx,x\rangle=0$.
\end{proof}

\section{Metric and orthogonal projections}

\begin{notedefinition}[Distance from a set]
For a nonempty subset $C$ of a normed space,
\[
  \dist(x,C)=\inf_{c\in C}\|x-c\|.
\]
A set is convex if it contains every line segment joining two of its points.
\end{notedefinition}

\begin{notetheorem}[Projection theorem for closed convex sets]
Let $C$ be a nonempty closed convex subset of a Hilbert space $H$.  For every
$x\in H$ there is a unique $P_Cx\in C$ such that
\[
  \|x-P_Cx\|=\dist(x,C).
\]
\end{notetheorem}

\begin{proof}
Let $d=\dist(x,C)$ and choose $c_n\in C$ with $\|x-c_n\|\to d$.  The
parallelogram identity gives
\[
 \|c_n-c_m\|^2
 =2\|x-c_n\|^2+2\|x-c_m\|^2
  -4\left\|x-\frac{c_n+c_m}{2}\right\|^2.
\]
The midpoint belongs to $C$, so the last norm is at least $d$; hence $(c_n)$
is Cauchy.  Completeness and closedness give a minimizer.  If $c,c'$ are two
minimizers, the same identity applied to them gives $\|c-c'\|=0$.
\end{proof}

\begin{notedefinition}[Orthogonal projection]
If $M$ is a closed subspace of a Hilbert space, the nearest-point map
$P_M\colon H\to M$ is the orthogonal projection onto $M$.
\end{notedefinition}

\begin{notetheorem}[Characterization of orthogonal projection]
Let $M$ be a closed subspace and $x\in H$.  For $m\in M$, the following are
equivalent:
\begin{enumerate}
  \item $m=P_Mx$;
  \item $x-m\perp M$.
\end{enumerate}
Moreover, $P_M$ is linear,
\[
  \|P_Mx\|\le\|x\|,
\]
and equality holds if and only if $x\in M$.
\end{notetheorem}

\begin{proof}
If $m$ is minimizing, then for every $y\in M$ and scalar $t$,
\[
  \|x-m\|^2\le\|x-m-ty\|^2.
\]
Varying real and imaginary $t$ forces $\langle x-m,y\rangle=0$.
Conversely, Pythagoras gives
$\|x-y\|^2=\|x-m\|^2+\|m-y\|^2$ for $y\in M$.
The orthogonality characterization proves linearity.  Finally,
$\|x\|^2=\|P_Mx\|^2+\|x-P_Mx\|^2$.
\end{proof}

\begin{notetheorem}[Variational inequality for convex projection]
Let $C$ be nonempty, closed, and convex.  For $x\in H$ and $y\in C$,
\[
  y=P_Cx
  \quad\Longleftrightarrow\quad
  \operatorname{Re}\langle x-y,z-y\rangle\le0
  \quad(z\in C).
\]
The projection is nonexpansive:
\[
  \|P_Cx-P_Cu\|\le\|x-u\|.
\]
It is linear if and only if $C$ is a closed subspace.
\end{notetheorem}

\begin{proof}
For $z\in C$, minimality against $y+t(z-y)$ and division by $t>0$ gives the
variational inequality as $t\downarrow0$.  Conversely, expand
\[
 \|x-z\|^2=\|x-y\|^2+\|z-y\|^2
            -2\operatorname{Re}\langle x-y,z-y\rangle.
\]
For $y=P_Cx$ and $v=P_Cu$, apply the inequality twice to obtain
\[
  \|y-v\|^2\le\operatorname{Re}\langle x-u,y-v\rangle
  \le\|x-u\|\|y-v\|.
\]
If $C$ is a subspace, this is the orthogonal projection and is linear.  The
range of any linear projection is a subspace, giving the converse.
\end{proof}

\begin{noteexercise}[Resolvent of a projection]
Let $P$ be a bounded projection, $P^2=P$.  If $\lambda\notin\{0,1\}$, then
$\lambda I-P$ is invertible and
\[
  (\lambda I-P)^{-1}
  =\frac1\lambda(I-P)+\frac1{\lambda-1}P.
\]
\end{noteexercise}

\begin{proof}
Decompose $x=(I-P)x+Px$ into $\Ker P\oplus\Ran P$.  On these two subspaces,
$\lambda I-P$ acts as multiplication by $\lambda$ and $\lambda-1$,
respectively.
\end{proof}

\begin{noteexercise}[Nested convex projections]
Let $C_1\supseteq C_2\supseteq\cdots$ be nonempty closed convex subsets of a
Hilbert space, with nonempty intersection $C$.  If $x_n=P_{C_n}x$, then
$x_n\to P_Cx$.
\end{noteexercise}

\begin{proof}
The distances $\|x-x_n\|$ increase and are bounded above by
$\dist(x,C)$.  The projection inequalities show $(x_n)$ is Cauchy.  Its limit
lies in every $C_n$ and attains the distance to $C$.
\end{proof}

\begin{noteremark}
Source pages 136--137 are blank and contain no mathematical material.
\end{noteremark}

\section{Orthogonal complements and Riesz representation}

\begin{notedefinition}[Orthogonal complement]
For $M\subseteq H$, define
\[
  M^\perp=\{x\in H:\langle x,m\rangle=0\text{ for every }m\in M\}.
\]
\end{notedefinition}

\begin{notetheorem}[Basic properties]
For subsets and subspaces of a Hilbert space:
\begin{enumerate}
  \item $M^\perp$ is a closed subspace;
  \item $M\cap M^\perp=\{0\}$ when $M$ is a subspace;
  \item $M\subseteq N$ implies $N^\perp\subseteq M^\perp$;
  \item $(\overline M)^\perp=M^\perp$;
  \item $M\subseteq(M^\perp)^\perp$.
\end{enumerate}
\end{notetheorem}

\begin{proof}
Linearity is immediate.  Closedness follows by passing inner products to the
limit.  The remaining statements follow directly from the definition.
\end{proof}

\begin{notetheorem}[Double orthogonal complement]
For every subspace $M$ of a Hilbert space,
\[
  (M^\perp)^\perp=\overline M.
\]
Thus $M$ is dense if and only if $M^\perp=\{0\}$.
\end{notetheorem}

\begin{proof}
The inclusion $\overline M\subseteq(M^\perp)^\perp$ is automatic.  If
$x\in(M^\perp)^\perp$, project $x$ onto $\overline M$.  The remainder
$x-P_{\overline M}x$ belongs both to $M^\perp$ and to its orthogonal
complement, hence is zero.
\end{proof}

\begin{notetheorem}[Orthogonal decomposition]
If $M$ is a closed subspace of a Hilbert space $H$, then
\[
  H=M\oplus M^\perp.
\]
Every $x$ has the unique form
\[
  x=P_Mx+(I-P_M)x,
\]
with the two summands in $M$ and $M^\perp$.
\end{notetheorem}

\begin{notetheorem}[Range and kernel of an orthogonal projection]
For a closed subspace $M$,
\[
  \Ran P_M=M,
  \qquad
  \Ker P_M=M^\perp,
  \qquad
  P_{M^\perp}=I-P_M.
\]
\end{notetheorem}

\begin{notetheorem}[Riesz representation theorem]
For every bounded linear functional $\varphi$ on a Hilbert space $H$, there
is a unique $h\in H$ such that
\[
  \varphi(x)=\langle x,h\rangle
  \qquad(x\in H).
\]
Moreover,
\[
  \|\varphi\|=\|h\|.
\]
\end{notetheorem}

\begin{proof}
If $\varphi=0$, take $h=0$.  Otherwise, $N=\Ker\varphi$ is a proper closed
subspace.  Choose a unit vector $u\in N^\perp$ and set
$h=\overline{\varphi(u)}u$; this conjugate is required by the convention that
the inner product is linear in the first variable.  Every $x$ decomposes as
$n+\alpha u$, and then
\[
  \langle x,h\rangle
  =\alpha\varphi(u)=\varphi(x).
\]
Uniqueness follows by testing the difference against itself.  Cauchy--Schwarz
gives $\|\varphi\|\le\|h\|$, while testing $h/\|h\|$ gives equality.
\end{proof}

\begin{notetheorem}[Representation of bounded sesquilinear forms]
Let $H_1,H_2$ be Hilbert spaces and let
$b\colon H_1\times H_2\to\mathbb K$ be bounded, linear in the first variable
and conjugate linear in the second.  Then there is a unique
$S\in\Bop{H_1}{H_2}$ such that
\[
  b(x,y)=\langle Sx,y\rangle,
  \qquad
  \|S\|=\|b\|.
\]
\end{notetheorem}

\begin{proof}
For each $x$, apply Riesz representation to the linear functional
$y\mapsto\overline{b(x,y)}$.  It gives a unique $Sx\in H_2$ such that
$\overline{b(x,y)}=\langle y,Sx\rangle$, equivalently
$b(x,y)=\langle Sx,y\rangle$.  Uniqueness gives linearity in $x$, and the
norm estimate gives boundedness and equality of norms.
\end{proof}

\begin{notecorollary}[Hilbert direct sum]
For every closed subspace $M$,
\[
  H=M\oplus M^\perp.
\]
\end{notecorollary}

\begin{noteexercise}[Strict convexity of the Hilbert ball]
The closed unit ball of an inner-product space is strictly convex: if
$x,y\in S_H$ and $x\ne y$, then
\[
  \left\|\frac{x+y}{2}\right\|<1.
\]
This can fail in a general normed space; for example in $\ell^\infty_2$ the
unit sphere contains line segments.
\end{noteexercise}

\begin{proof}
The parallelogram identity gives
\[
  \left\|\frac{x+y}{2}\right\|^2
  =1-\frac14\|x-y\|^2.
\]
\end{proof}

\begin{notetheorem}[Hilbert-space extension of operators]
Let $U$ be a closed subspace of a Hilbert space $H$, let $W$ be Banach, and
let $S\in\Bop{U}{W}$.  Then
\[
  T=S\circ P_U
\]
is a bounded extension of $S$ to $H$ and
\[
  \|T\|=\|S\|.
\]
For a nonclosed $U$, first extend $S$ continuously to $\overline U$.
\end{notetheorem}

\begin{proof}
Since $P_U|_U=I$ and $\|P_U\|\le1$, one has
$\|T\|\le\|S\|$.  Restriction to $U$ gives the reverse inequality.
\end{proof}

\begin{noteexercise}[Equality in the triangle inequality for functionals]
Let $\varphi,\psi\in H^*$ be nonzero and suppose
\[
  \|\varphi+\psi\|=\|\varphi\|+\|\psi\|.
\]
Then one is a nonnegative real multiple of the other.  The conclusion need not
hold in an arbitrary Banach space.
\end{noteexercise}

\begin{proof}
Represent the functionals by $h,k\in H$.  The hypothesis becomes equality in
the Hilbert-space triangle inequality for $h+k$, so $h$ and $k$ are
nonnegative real multiples.  In $\ell^\infty_2{}^*=\ell^1_2$, equality can
hold for independent positive coordinate functionals.
\end{proof}

\section{Orthonormal families and bases}

\begin{notedefinition}[Orthonormal family]
A family $(e_i)_{i\in I}$ is orthonormal if
\[
  \langle e_i,e_j\rangle=\delta_{ij}.
\]
For finite scalar families,
\[
  \left\|\sum_{i\in F}a_ie_i\right\|^2
  =\sum_{i\in F}|a_i|^2.
\]
\end{notedefinition}

\begin{notedefinition}[Unordered sum]
For a family $(x_i)_{i\in I}$ in a normed space, the unordered sum
$\sum_{i\in I}x_i$ converges to $x$ if for every $\varepsilon>0$ there is a
finite $F_0\subseteq I$ such that
\[
  \left\|x-\sum_{i\in F}x_i\right\|<\varepsilon
\]
for every finite $F\supseteq F_0$.
\end{notedefinition}

\begin{notetheorem}[Linear combinations of an orthonormal family]
Let $(e_i)_{i\in I}$ be orthonormal in a Hilbert space.  The unordered series
\[
  \sum_{i\in I}a_ie_i
\]
converges if and only if
\[
  \sum_{i\in I}|a_i|^2<\infty.
\]
When it converges,
\[
  \left\|\sum_{i\in I}a_ie_i\right\|^2
  =\sum_{i\in I}|a_i|^2.
\]
\end{notetheorem}

\begin{proof}
Differences of finite partial sums have squared norm equal to the corresponding
finite sum of $|a_i|^2$.  Thus the net of finite partial sums is Cauchy exactly
when the scalar unordered sum is finite.  Completeness gives convergence and
passing to the limit gives the norm identity.  Only countably many $a_i$ can
be nonzero.
\end{proof}

\begin{notetheorem}[Bessel's inequality]
For an orthonormal family $(e_i)_{i\in I}$ and $x\in H$,
\[
  \sum_{i\in I}|\langle x,e_i\rangle|^2\le\|x\|^2.
\]
\end{notetheorem}

\begin{proof}
For a finite $F$, let
$y_F=\sum_{i\in F}\langle x,e_i\rangle e_i$.  Then
$x-y_F\perp y_F$, so
\[
  \|x\|^2=\|x-y_F\|^2+\sum_{i\in F}|\langle x,e_i\rangle|^2.
\]
Take the supremum over finite $F$.
\end{proof}

\begin{notetheorem}[Closed span]
The closed span of an orthonormal family is
\[
  \overline{\Span\{e_i:i\in I\}}
  =\left\{\sum_{i\in I}a_ie_i:
       \sum_i|a_i|^2<\infty\right\}.
\]
For $x$ in this closed span,
\[
  x=\sum_{i\in I}\langle x,e_i\rangle e_i.
\]
\end{notetheorem}

\begin{notetheorem}[Maximal orthonormal families]
An orthonormal family $(e_i)$ is maximal if and only if its closed linear span
is all of $H$.
\end{notetheorem}

\begin{proof}
If the closed span is proper, choose a nonzero vector in its orthogonal
complement and normalize it; adjoining that vector enlarges the family.
Conversely, a vector orthogonal to a dense span is zero.
\end{proof}

\begin{notedefinition}[Orthonormal basis]
An orthonormal family with dense linear span is an orthonormal basis.
\end{notedefinition}

\begin{notetheorem}[Parseval identities]
If $(e_i)_{i\in I}$ is an orthonormal basis, then for $x,y\in H$,
\[
  x=\sum_{i\in I}\langle x,e_i\rangle e_i,
\]
\[
  \langle x,y\rangle
  =\sum_{i\in I}\langle x,e_i\rangle
      \overline{\langle y,e_i\rangle},
\]
and
\[
  \|x\|^2=\sum_{i\in I}|\langle x,e_i\rangle|^2.
\]
\end{notetheorem}

\begin{notetheorem}[Totality]
A subset $M\subseteq H$ is total if and only if $M^\perp=\{0\}$.  If $H$ is
complete, this is equivalent to $\overline{\Span M}=H$.
\end{notetheorem}

\begin{notetheorem}[Orthonormal families in separable spaces]
Every orthonormal family in a separable Hilbert space is countable.  Every
separable Hilbert space has a countable orthonormal basis.
\end{notetheorem}

\begin{proof}
Distinct orthonormal vectors have distance $\sqrt2$.  Balls of radius
$1/3$ around them are disjoint, and each contains a distinct member of a
fixed countable dense set.  Thus the family is countable.  For existence,
apply Gram--Schmidt to a dense sequence after discarding vectors already in
the span of their predecessors.
\end{proof}

\begin{notetheorem}[Classification by Hilbert dimension]
Two Hilbert spaces are isometrically isomorphic if and only if their
orthonormal bases have the same cardinality.
\end{notetheorem}

\begin{proof}
First, any two orthonormal bases of one Hilbert space have the same
cardinality.  Indeed, if $E$ and $F$ are such bases, then for each $e\in E$
the set $\{f\in F:\langle e,f\rangle\ne0\}$ is countable by Bessel's
inequality.  Totality of $E$ shows that the union of these countable sets is
all of $F$, so $|F|\le |E|$ for infinite bases; reversing the roles gives
equality, while the finite-dimensional case is elementary.

Given equally indexed bases $(e_i)_{i\in I}$ and $(f_i)_{i\in I}$, define
\[
  T\!\left(\sum_ia_ie_i\right)=\sum_ia_if_i.
\]
Parseval shows that $T$ is a surjective linear isometry.  Conversely, a
surjective linear isometry sends an orthonormal basis to an orthonormal basis,
so the cardinalities agree.
\end{proof}

\begin{notetheorem}[Projection and Fourier coefficients]
Let $(e_i)_{i\in I}$ be an orthonormal basis of a closed subspace $M$.  Then
\[
  P_Mx=\sum_{i\in I}\langle x,e_i\rangle e_i.
\]
If $\varphi\in H^*$ and $\varphi(e_i)=a_i$, then
$(a_i)\in\ell^2(I)$ and the Riesz representative is
\[
  h=\sum_{i\in I}\overline{a_i}e_i,
\]
with $\|\varphi\|^2=\sum_i|a_i|^2$ under the linear-first convention.
\end{notetheorem}

\begin{proof}
The difference between $x$ and the displayed sum is orthogonal to every basis
vector and hence to $M$.  The functional statement follows by applying the
Riesz theorem and Parseval to its representing vector.
\end{proof}

\begin{noteexercise}[No nonzero vector orthogonal to a total family]
If $M$ is total and $x\perp M$, then $x=0$.
\end{noteexercise}

\begin{proof}
Orthogonality to $M$ implies orthogonality to its closed span, which is all of
$H$; in particular $\langle x,x\rangle=0$.
\end{proof}

\chapter{Spectral Theory on Normed Spaces}

Throughout the spectral portions of this chapter, scalar fields are complex
unless explicitly stated otherwise.  For a linear operator
$T\colon\Dom(T)\subseteq X\to X$ and $\lambda\in\mathbb C$, write
\[
  T_\lambda:=T-\lambda I.
\]
When $T_\lambda$ is one-to-one, its inverse is understood as a map from
$\Ran(T_\lambda)$ to $\Dom(T)$.

\section{Resolvents and spectra}

\begin{notedefinition}[Regular values, resolvent, and spectrum]
Let $X\ne\{0\}$ be a complex Banach space and let
$T\colon\Dom(T)\subseteq X\to X$ be a closed linear operator.  A number
$\lambda\in\mathbb C$ is a \emph{regular value} if
$T_\lambda\colon\Dom(T)\to X$ is bijective and its inverse
\[
  R_\lambda(T):=T_\lambda^{-1}\colon X\to X
\]
is bounded.  The regular values form the \emph{resolvent set} $\rho(T)$,
and
\[
  \sigma(T):=\mathbb C\setminus\rho(T)
\]
is the \emph{spectrum}.

The spectrum is divided as follows:
\begin{itemize}
  \item $\lambda$ lies in the point spectrum $\sigma_p(T)$ when
        $T_\lambda$ is not injective; these are precisely the eigenvalues;
  \item it lies in the continuous spectrum when $T_\lambda$ is injective
        and has dense, non-surjective range;
  \item it lies in the residual spectrum when $T_\lambda$ is injective but
        its range is not dense.
\end{itemize}
The eigenspace corresponding to $\lambda$ is $\Ker(T_\lambda)$.
\end{notedefinition}

\begin{noteremark}[The convention used in the source]
The source calls $\lambda$ regular when $T_\lambda$ is injective and its
inverse is bounded on a dense range.  For a closed operator on a Banach space,
that convention is equivalent to the standard one above, as the next lemma
shows.  Bounded everywhere-defined operators are closed, so no distinction
arises in the remainder of these notes.
\end{noteremark}

\begin{notelemma}[A bounded inverse on a dense range]
Let $X$ be a Banach space, let
$T\colon\Dom(T)\subseteq X\to X$ be closed, and suppose $T_\lambda$ is
injective with dense range.  If
$T_\lambda^{-1}\colon\Ran(T_\lambda)\to X$ is bounded, then
$\Ran(T_\lambda)=X$.
\end{notelemma}

\begin{proof}
The inverse is closed because its graph is obtained from
$\Graph(T_\lambda)$ by interchanging the coordinates.  Its domain is also
closed: if $y_n\in\Ran(T_\lambda)$ and $y_n\to y$, boundedness makes
$T_\lambda^{-1}y_n$ Cauchy, hence convergent in $X$, and closedness of the
inverse puts $y$ back in its domain.  The range is therefore both dense and
closed, so it is all of $X$.
\end{proof}

\begin{noteexample}[A diagonal operator on $\ell^2$]
Let $(\alpha_j)$ be dense in $[0,1]$ and define
\[
  T(\xi_j)=(\alpha_j\xi_j)
  \qquad\text{on }\ell^2.
\]
Then
\[
  \sigma_p(T)=\{\alpha_j:j\ge1\},
  \qquad
  \sigma(T)=\overline{\{\alpha_j:j\ge1\}}=[0,1].
\]
If $\lambda\in[0,1]\setminus\{\alpha_j\}$, choose $j_n$ with
$\alpha_{j_n}\to\lambda$.  Then
\[
  R_\lambda(T)e_{j_n}
  =\frac1{\alpha_{j_n}-\lambda}e_{j_n},
\]
so the inverse is unbounded and $\lambda$ lies in the continuous spectrum.
\end{noteexample}

\begin{notetheorem}[Neumann series]
If $T\in\Bop{X}{X}$ on a Banach space and $\|T\|<1$, then
\[
  (I-T)^{-1}=\sum_{n=0}^{\infty}T^n
\]
with convergence in operator norm.
\end{notetheorem}

\begin{proof}
The series converges absolutely in the Banach algebra $\Bop{X}{X}$.  Its
partial sums satisfy
\[
  (I-T)\sum_{j=0}^{N}T^j=I-T^{N+1}
  =\left(\sum_{j=0}^{N}T^j\right)(I-T),
\]
and $T^{N+1}\to0$ in norm.
\end{proof}

\begin{notetheorem}[Openness of the resolvent and local expansion]
Let $T\in\Bop{X}{X}$ on a complex Banach space.  The resolvent set is open.
If $\lambda_0\in\rho(T)$ and
\[
  |\lambda-\lambda_0|<\frac1{\|R_{\lambda_0}(T)\|},
\]
then $\lambda\in\rho(T)$ and
\[
  R_\lambda(T)
  =\sum_{j=0}^{\infty}
     (\lambda-\lambda_0)^jR_{\lambda_0}(T)^{j+1}.
\]
Consequently $\sigma(T)$ is closed and the resolvent is operator-valued
holomorphic on $\rho(T)$.
\end{notetheorem}

\begin{proof}
Factor
\[
  T-\lambda I
  =(T-\lambda_0I)
   \bigl[I-(\lambda-\lambda_0)R_{\lambda_0}(T)\bigr]
\]
and apply the Neumann series to the second factor.
\end{proof}

\begin{notedefinition}[Spectral radius]
For $T\in\Bop{X}{X}$, the spectral radius is
\[
  r(T):=\sup\{ |\lambda|:\lambda\in\sigma(T)\}.
\]
\end{notedefinition}

\begin{notetheorem}[Resolvent identity and commutation]
For $\lambda,\mu\in\rho(T)$,
\[
  R_\mu(T)-R_\lambda(T)
  =(\mu-\lambda)R_\mu(T)R_\lambda(T).
\]
If $S\in\Bop{X}{X}$ commutes with $T$, then it commutes with every
$R_\lambda(T)$.  In particular, all resolvents commute with one another.
\end{notetheorem}

\begin{proof}
The first formula follows from
\[
  A^{-1}-B^{-1}=A^{-1}(B-A)B^{-1}
\]
with $A=T-\mu I$ and $B=T-\lambda I$.  If $ST=TS$, then
$S(T-\lambda I)=(T-\lambda I)S$; multiplying by the inverse on both sides
gives $SR_\lambda=R_\lambda S$.
\end{proof}

\begin{notetheorem}[Polynomial spectral mapping]
For a polynomial $p$ and $T\in\Bop{X}{X}$,
\[
  \sigma(p(T))=p(\sigma(T)).
\]
\end{notetheorem}

\begin{proof}
If $\mu\notin p(\sigma(T))$, factor
\[
  p(z)-\mu=c\prod_{j=1}^{m}(z-\lambda_j).
\]
Every $\lambda_j$ belongs to $\rho(T)$, so the commuting product
$p(T)-\mu I$ is invertible.

Conversely, if $\mu=p(\lambda)$ with $\lambda\in\sigma(T)$, polynomial
division gives
\[
  p(T)-p(\lambda)I=(T-\lambda I)q(T).
\]
If the left-hand side were invertible, the commuting factors would give both
a left and a right inverse for $T-\lambda I$, contradicting
$\lambda\in\sigma(T)$.
\end{proof}

\begin{notelemma}[Eigenvectors for distinct eigenvalues]
Eigenvectors corresponding to distinct eigenvalues are linearly independent.
\end{notelemma}

\begin{noteexample}[Idempotent operators]
If $T^2=T$ and $T\ne0,I$, then
\[
  \sigma(T)=\{0,1\}.
\]
Indeed, polynomial spectral mapping applied to $p(z)=z^2-z$ gives
$\sigma(T)\subseteq\{0,1\}$.  Since $T$ and $T-I$ are both noninjective,
$0$ and $1$ are eigenvalues.
\end{noteexample}

\begin{notetheorem}[Resolvent norm and distance to the spectrum]
For $\lambda\in\rho(T)$,
\[
  \|R_\lambda(T)\|
  \ge \frac1{\dist(\lambda,\sigma(T))}.
\]
Hence the resolvent norm tends to infinity as $\lambda$ approaches the
spectrum through $\rho(T)$.
\end{notetheorem}

\begin{proof}
The local expansion shows that the open disk of radius
$1/\|R_\lambda(T)\|$ centred at $\lambda$ lies in $\rho(T)$.  Its radius
cannot exceed the distance from $\lambda$ to the complementary set
$\sigma(T)$.
\end{proof}

\begin{notetheorem}[Nonempty compact spectrum]
If $X\ne\{0\}$ is a complex Banach space and $T\in\Bop{X}{X}$, then
$\sigma(T)$ is nonempty and compact, and
\[
  \sigma(T)\subseteq\{\lambda:|\lambda|\le\|T\|\}.
\]
\end{notetheorem}

\begin{proof}
For $|\lambda|>\|T\|$ the Neumann expansion
\[
  R_\lambda(T)
  =-\frac1\lambda\sum_{n=0}^{\infty}\frac{T^n}{\lambda^n}
\]
exists, so the spectrum is bounded.  It is closed by openness of the
resolvent.

If the spectrum were empty, then for every $x\in X$ and $f\in X^*$ the
scalar function $h(\lambda)=f(R_\lambda(T)x)$ would be entire.  The displayed
expansion shows that $h(\lambda)\to0$ as $|\lambda|\to\infty$, so Liouville's
theorem makes $h$ identically zero.  Since bounded functionals separate
points, $R_\lambda(T)=0$, impossible for an inverse.
\end{proof}

\begin{notetheorem}[Spectral radius formula]
For $T\in\Bop{X}{X}$ on a complex Banach space,
\[
  r(T)=\lim_{n\to\infty}\|T^n\|^{1/n}
      =\inf_{n\ge1}\|T^n\|^{1/n}.
\]
\end{notetheorem}

\begin{proof}
The sequence $a_n=\|T^n\|$ is submultiplicative, so
$a_n^{1/n}$ has a limit equal to its infimum.  Polynomial spectral mapping
gives
\[
  r(T)^n=r(T^n)\le\|T^n\|,
\]
therefore $r(T)\le\lim_n\|T^n\|^{1/n}$.

For the reverse inequality, fix $\rho>r(T)$.  The circle
$|\lambda|=\rho$ lies in the resolvent set.  The Laurent coefficients of the
resolvent at infinity satisfy the Cauchy formula
\[
  T^n=-\frac1{2\pi i}\int_{|\lambda|=\rho}
       \lambda^nR_\lambda(T)\,\mathrm d\lambda.
\]
Thus, with $M_\rho=\max_{|\lambda|=\rho}\|R_\lambda(T)\|$,
\[
  \|T^n\|\le M_\rho\rho^{n+1}.
\]
Taking $n$th roots and then letting $\rho\downarrow r(T)$ gives the reverse
inequality.
\end{proof}

\section{Compact operators}

\begin{notedefinition}[Compact linear operator]
Let $X,Y$ be normed spaces.  A linear operator $T\colon X\to Y$ is
\emph{compact} if it maps every bounded subset of $X$ to a relatively compact
subset of $Y$.
\end{notedefinition}

\begin{notetheorem}[Basic compactness criteria]
Let $T\colon X\to Y$ be linear.
\begin{enumerate}
  \item $T$ is compact if and only if every bounded sequence $(x_n)$ has a
        subsequence $(x_{n_k})$ for which $(Tx_{n_k})$ converges in $Y$.
  \item Every compact linear operator is bounded.
  \item The identity on an infinite-dimensional normed space is not compact.
  \item If $T$ is bounded and either its domain or its range is
        finite-dimensional, then $T$ is compact.
\end{enumerate}
\end{notetheorem}

\begin{proof}
The first assertion is sequential compactness in metric spaces.  For the
second, the relatively compact set $T(B_X)$ is bounded, so
$\sup_{\|x\|\le1}\|Tx\|<\infty$.  The third follows because the closed unit
ball is not compact in infinite dimension.  The final assertion follows from
finite-dimensional Heine--Borel compactness.
\end{proof}

\begin{notetheorem}[Norm limits of compact operators]
Let $Y$ be Banach.  If $T_n\in\Kop{X}{Y}$ and $\|T_n-T\|\to0$, then $T$ is
compact.  Consequently $\Kop{X}{Y}$ is a closed subspace of
$\Bop{X}{Y}$.
\end{notetheorem}

\begin{proof}
Given a bounded sequence, successively extract subsequences on which
$T_1,T_2,\ldots$ converge, and take the diagonal subsequence $(x_{n_k})$.
Put $M=\max\{1,\sup_k\|x_{n_k}\|\}$.  Choose $j$ with
$\|T-T_j\|<\varepsilon/(3M)$ and then $k,\ell$ large enough that
$\|T_jx_{n_k}-T_jx_{n_\ell}\|<\varepsilon/3$.  It follows that
$(Tx_{n_k})$ is Cauchy, hence convergent in $Y$.
\end{proof}

\begin{notetheorem}[Weak-to-strong continuity]
If $T\colon X\to Y$ is compact and $x_n\weakto x$, then
\[
  Tx_n\longrightarrow Tx
\]
in norm.
\end{notetheorem}

\begin{proof}
The sequence $(x_n)$ is bounded.  Every subsequence of $(Tx_n)$ has a
norm-convergent further subsequence by compactness.  Its norm limit must be
$Tx$, because bounded linear maps preserve weak convergence and weak limits
are unique.  Therefore the whole sequence converges to $Tx$ in norm.
\end{proof}

\begin{notedefinition}[$\varepsilon$-nets and total boundedness]
A finite set $F$ is an $\varepsilon$-net for $B$ if every point of $B$ lies
within distance $\varepsilon$ of a point of $F$.  A set is \emph{totally
bounded} if it has a finite $\varepsilon$-net for every $\varepsilon>0$.
\end{notedefinition}

\begin{notelemma}[Total boundedness]
For a subset $B$ of a metric space:
\begin{enumerate}
  \item relative compactness implies total boundedness;
  \item total boundedness in a complete space implies relative compactness;
  \item the centres of a finite $\varepsilon$-net may be chosen in $B$;
  \item every totally bounded set is separable.
\end{enumerate}
\end{notelemma}

\begin{proof}
A compact set is covered by finitely many $\varepsilon$-balls.  Conversely,
successive finite $2^{-n}$-nets yield a Cauchy subsequence of every sequence;
completeness supplies its limit.  Replacing each nonempty covering ball by a
point of $B$ proves the third assertion.  The union of finite $1/n$-nets is a
countable dense subset.
\end{proof}

\begin{notetheorem}[Separable range and extension to the completion]
Let $T\colon X\to Y$ be compact.
\begin{enumerate}
  \item $\Ran(T)$ is separable.
  \item If $Y$ is Banach and $\widehat X$ is the completion of $X$, the
        bounded extension $\widehat T\colon\widehat X\to Y$ is compact.
\end{enumerate}
\end{notetheorem}

\begin{proof}
Since $X=\bigcup_{n\ge1}nB_X$ and each $T(nB_X)$ is relatively compact and
therefore separable, a countable union of countable dense sets is dense in
$\Ran(T)$.

For the extension, let $(\widehat x_n)$ be bounded in $\widehat X$ and choose
$x_n\in X$ with $\|x_n-\widehat x_n\|\to0$.  The sequence $(x_n)$ is bounded,
so some $Tx_{n_k}$ converges.  Since
$\widehat T(\widehat x_{n_k}-x_{n_k})\to0$, the corresponding
$\widehat T\widehat x_{n_k}$ also converges.
\end{proof}

\begin{notetheorem}[Schauder theorem]
If $T\colon X\to Y$ is compact, then its adjoint
$T^*\colon Y^*\to X^*$ is compact.
\end{notetheorem}

\begin{proof}
The set $K=\overline{T(B_X)}$ is compact and hence totally bounded.  Given
$\varepsilon>0$, choose $y_1,\ldots,y_N\in K$ so that every $Tx$ with
$\|x\|\le1$ is within $\varepsilon/4$ of some $y_j$.  On the dual unit ball,
the coordinate map
\[
  g\longmapsto(g(y_1),\ldots,g(y_N))
\]
has totally bounded image in the finite-dimensional space $\mathbb K^N$.
Choose finitely many representatives so that any two functionals in the same
coordinate cell differ by less than $\varepsilon/2$ on each $y_j$.
Then, for $\|x\|\le1$ and the corresponding $j$,
\[
 |(g-h)(Tx)|
 \le |(g-h)(y_j)|+\|g-h\|\,\|Tx-y_j\|<\varepsilon.
\]
Thus $T^*(B_{Y^*})$ is totally bounded in the Banach space $X^*$, hence
relatively compact.
\end{proof}

\section{Spectral structure of compact operators}

\begin{notetheorem}[Nonzero eigenvalues]
Let $T\colon X\to X$ be compact.  Its set of eigenvalues is at most
countable, and $0$ is its only possible accumulation point.  Equivalently,
for every $c>0$ there are only finitely many eigenvalues satisfying
$|\lambda|\ge c$.
\end{notetheorem}

\begin{proof}
Suppose distinct eigenvalues $\lambda_n$ satisfy $|\lambda_n|\ge c>0$, and
choose eigenvectors so that
$M_n=\Span\{x_1,\ldots,x_n\}$ is strictly increasing.  By Riesz's lemma choose
$y_n\in M_n$ with $\|y_n\|=1$ and
$\dist(y_n,M_{n-1})\ge1/2$.  Each $M_n$ is $T$-invariant.  For $m<n$,
\[
  Ty_n-Ty_m
  =\lambda_n y_n+z_{m,n}
  \qquad(z_{m,n}\in M_{n-1}),
\]
so $\|Ty_n-Ty_m\|\ge c/2$.  This contradicts compactness.  The finiteness
claim implies countability by taking $c=1/k$.
\end{proof}

\begin{notelemma}[Composition]
If $T$ is compact and $S$ is bounded, then $TS$ and $ST$ are compact whenever
the compositions are defined.
\end{notelemma}

\begin{notetheorem}[Finite-dimensional null spaces]
Let $T\colon X\to X$ be compact and $\lambda\ne0$.  Then
\[
  \dim\Ker(T-\lambda I)<\infty.
\]
More generally, every $\Ker(T-\lambda I)^n$ is finite-dimensional.
\end{notetheorem}

\begin{proof}
On $N=\Ker(T-\lambda I)$ one has $Tx=\lambda x$.  Hence the image under $T$
of the closed unit ball of $N$ is a scalar multiple of that ball.  Compactness
of $T$ makes the ball compact, so $N$ is finite-dimensional.

For powers, the binomial expansion gives
\[
  (T-\lambda I)^n=(-\lambda)^nI+K_n,
\]
where $K_n$ is compact.  The same argument applied to $K_n$ shows that the
kernel is finite-dimensional.
\end{proof}

\begin{notetheorem}[Closed range of a compact perturbation of a scalar]
Let $X$ be Banach, let $T\colon X\to X$ be compact, and let
$\lambda\ne0$.  Then
\[
  \Ran(T-\lambda I)
\]
is closed.  The same is true for every power $(T-\lambda I)^n$.
\end{notetheorem}

\begin{proof}
Put $A=T-\lambda I$ and $N=\Ker A$.  The induced injective map
\[
  \widetilde A\colon X/N\to X,
  \qquad \widetilde A(x+N)=Ax,
\]
is bounded below.  Otherwise there are cosets of norm one represented by a
bounded sequence $(x_n)$ with $Ax_n\to0$.  Compactness gives a subsequence
for which $Tx_n$ converges, and then
\[
  \lambda x_n=Tx_n-Ax_n
\]
converges.  Its limit belongs to $N$, contradicting
$\dist(x_n,N)=1$.  A bounded-below operator has closed range.

For $A^n=(-\lambda)^nI+K_n$, apply the same argument to the compact operator
$K_n$.
\end{proof}

\begin{notecorollary}
Let $T\colon X\to Y$ be compact between Banach spaces.  Then
$\Ran(T)$ contains no closed infinite-dimensional subspace of $Y$.
\end{notecorollary}

\begin{proof}
If a closed infinite-dimensional subspace $M\subseteq\Ran(T)$ existed, the
map $T\colon T^{-1}(M)\to M$ would be a bounded surjection between Banach
spaces.  The open mapping theorem would make a fixed multiple of the image of
the unit ball contain $B_M$.  That would make $B_M$ relatively compact,
contradicting infinite dimensionality.
\end{proof}

\section{Stabilization and the Riesz--Schauder theorem}

Throughout this section, $X$ is a Banach space and
$A=T-\lambda I$, where $T\colon X\to X$ is compact and $\lambda\ne0$.
Write
\[
  N_n:=\Ker A^n,
  \qquad
  R_n:=\Ran A^n.
\]
Then $N_0\subseteq N_1\subseteq\cdots$ and
$X=R_0\supseteq R_1\supseteq\cdots$.

\begin{notelemma}[Stabilization of null spaces]
There is a least integer $p\ge0$ such that
\[
  N_p=N_{p+1}=N_{p+2}=\cdots.
\]
If $p>0$, all inclusions $N_0\subsetneq N_1\subsetneq\cdots\subsetneq N_p$
are proper.
\end{notelemma}

\begin{proof}
If the chain were strictly increasing forever, Riesz's lemma would provide
$x_n\in N_n$ with $\|x_n\|=1$ and
$\dist(x_n,N_{n-1})\ge1/2$.  For $m<n$,
$Ax_n\in N_{n-1}$ and $Tx_m=Ax_m+\lambda x_m\in N_m\subseteq N_{n-1}$.
Thus
\[
  Tx_n-Tx_m=\lambda x_n+z_{m,n}
  \qquad(z_{m,n}\in N_{n-1}),
\]
so $\|Tx_n-Tx_m\|\ge|\lambda|/2$, contradicting compactness.  Once
$N_p=N_{p+1}$, applying $A$ shows inductively that every later equality also
holds.
\end{proof}

\begin{notelemma}[Stabilization of ranges]
There is a least integer $q\ge0$ such that
\[
  R_q=R_{q+1}=R_{q+2}=\cdots.
\]
If $q>0$, all inclusions $R_0\supsetneq R_1\supsetneq\cdots\supsetneq R_q$
are proper.
\end{notelemma}

\begin{proof}
The adjoint $T^*$ is compact and
\[
  A^*=T^*-\lambda I.
\]
The null-space chain of $A^*$ stabilizes.  Since every $R_n$ is closed and
\[
  R_n^\perp=\Ker(A^*)^n,
\]
equality of two successive adjoint kernels is equivalent, by Hahn--Banach,
to equality of the corresponding ranges.
\end{proof}

\begin{notetheorem}[Equal ascent and descent]
The least stabilization indices above are equal: $p=q=:r$.  Consequently
\[
  \Ker A^r=\Ker A^{r+1}=\cdots,
  \qquad
  \Ran A^r=\Ran A^{r+1}=\cdots.
\]
If $r>0$, every preceding inclusion in both chains is proper.
\end{notetheorem}

\begin{proof}
If the null chain stabilizes at $p$, then
\[
  \Ker A^p\cap\Ran A^p=\{0\}.
\]
Indeed, if $x=A^py$ and $A^px=0$, then $A^{2p}y=0$; stabilization gives
$y\in\Ker A^p$, hence $x=0$.

If the range chain stabilizes at $q$, then
\[
  X=\Ker A^q+\Ran A^q.
\]
For $x\in X$, equality $\Ran A^q=\Ran A^{2q}$ gives
$A^qx=A^{2q}y$ for some $y$, and then
$x-A^qy\in\Ker A^q$.

If $p\le q$, then $\Ker A^p=\Ker A^q$.  Applying $A^p$ to
$X=\Ker A^q+\Ran A^q$ gives
$\Ran A^p=\Ran A^{p+q}=\Ran A^q$, so $q\le p$.  If $q\le p$, decompose every
$x\in\Ker A^p$ as $x=u+v$ with $u\in\Ker A^q$ and
$v\in\Ran A^q=\Ran A^p$.  Then
$v\in\Ker A^p\cap\Ran A^p=\{0\}$, so
$\Ker A^p=\Ker A^q$ and $p\le q$.  Hence $p=q$.
\end{proof}

\begin{notetheorem}[Nonzero spectral values are eigenvalues]
Let $X$ be a complex Banach space and $T\in\Kop{X}{X}$.  Every
$\lambda\in\sigma(T)\setminus\{0\}$ is an eigenvalue.
\end{notetheorem}

\begin{proof}
If $A=T-\lambda I$ were injective, its null-space chain would stabilize at
$r=0$.  Equal ascent and descent would then give
$\Ran A=X$.  Thus $A$ would be a bounded bijection of Banach spaces, and the
bounded inverse theorem would put $\lambda$ in the resolvent, a contradiction.
\end{proof}

\begin{notetheorem}[Riesz decomposition]
With $r$ as above,
\[
  X=\Ker A^r\oplus\Ran A^r.
\]
Both summands are closed and invariant under $A$ and $T$.
\end{notetheorem}

\begin{proof}
The intersection is zero by the argument in the preceding theorem, and the
sum is all of $X$ by range stabilization.  Closedness of the two summands was
proved above.
\end{proof}

\begin{noteexercise}[Two noncompact operators]
\begin{enumerate}
  \item For the diagonal operator on $\ell^2$ with $(\alpha_j)$ dense in
        $[0,1]$, choose a subsequence with $|\alpha_{j_k}|\ge1/2$.  Then
        $\|Te_{j_k}-Te_{j_\ell}\|\ge1/\sqrt2$, so $T$ is not compact.
  \item On $C[0,1]$, the multiplication operator $(Tx)(t)=tx(t)$ is not
        compact.  The bounded sequence $x_n(t)=t^n$ has images
        $Tx_n=t^{n+1}$.  Every subsequence converges pointwise to the
        discontinuous function that is $0$ on $[0,1)$ and $1$ at $1$, so no
        subsequence can converge uniformly.
\end{enumerate}
\end{noteexercise}

\section{Fredholm solvability}

Throughout this section, $X$ is a complex Banach space.  Let
$T\in\Kop{X}{X}$, let $\lambda\ne0$, and put
\[
  A:=T-\lambda I.
\]
For both real and complex Banach spaces,
\[
  A^*=T^*-\lambda I.
\]
There is no conjugation here: the Banach-space adjoint acts on complex-linear
functionals, so $(\lambda I)^*=\lambda I$.  Conjugation appears only after
identifying a Hilbert space with its dual through the Riesz map.

\begin{notetheorem}[Solvability of $Ax=y$]
The equation
\[
  Ax=y
\]
has a solution if and only if
\[
  f(y)=0
  \qquad\text{for every }f\in\Ker A^*.
\]
In other words,
\[
  \Ran A={}^{\perp}(\Ker A^*).
\]
\end{notetheorem}

\begin{proof}
The identity $(\Ran A)^\perp=\Ker A^*$ holds for every bounded operator.  Here
$\Ran A$ is closed, so the Hahn--Banach annihilator theorem gives
\[
  \Ran A={}^{\perp}\bigl((\Ran A)^\perp\bigr)
        ={}^{\perp}(\Ker A^*).
\]
\end{proof}

\begin{notedefinition}[Normally solvable equation]
A bounded operator equation $Bx=y$ is called \emph{normally solvable} when
\[
  y\in\Ran B
  \quad\Longleftrightarrow\quad
  f(y)=0\text{ for every }f\in\Ker B^*.
\]
Equivalently, $\Ran B$ is closed.
\end{notedefinition}

\begin{notelemma}[A bounded choice of solution]
There is $C>0$ such that whenever $y\in\Ran A$, one can choose a solution
$\widetilde x$ of $A\widetilde x=y$ satisfying
\[
  \|\widetilde x\|\le C\|y\|.
\]
\end{notelemma}

\begin{proof}
The finite-dimensional space $\Ker A$ has a closed complement $M$.  The
restriction
\[
  A|_M\colon M\to\Ran A
\]
is a bounded bijection between Banach spaces.  Its inverse is bounded, and
$\widetilde x=(A|_M)^{-1}y$ has the required estimate.
\end{proof}

\begin{notetheorem}[Solvability of the adjoint equation]
For $g\in X^*$, the equation
\[
  A^*f=g
\]
has a solution if and only if
\[
  g(x)=0
  \qquad\text{for every }x\in\Ker A.
\]
Equivalently,
\[
  \Ran A^*=(\Ker A)^\perp.
\]
\end{notetheorem}

\begin{proof}
Define on $\Ran A$ the functional $F(Ax)=g(x)$.  The condition on $g$
makes $F$ well defined.  For $y=Ax$, choose a solution $\widetilde x$ with
$\|\widetilde x\|\le C\|y\|$; then
\[
  |F(y)|=|g(\widetilde x)|\le C\|g\|\|y\|.
\]
Thus $F$ is bounded.  Extend it to $f\in X^*$ by Hahn--Banach.  For every
$x\in X$,
$(A^*f)(x)=f(Ax)=F(Ax)=g(x)$, so $A^*f=g$.
\end{proof}

\begin{notetheorem}[Injectivity, surjectivity, and bounded inverses]
The following are equivalent:
\begin{enumerate}
  \item $A$ is injective;
  \item $A$ is surjective;
  \item $A$ is bijective.
\end{enumerate}
When these conditions hold, $A^{-1}$ is bounded.  The same equivalence holds
for $A^*$.
\end{notetheorem}

\begin{proof}
If $A$ is injective, its ascent is zero; equal ascent and descent imply that
its range chain stabilizes at zero, so $\Ran A=X$.  Conversely, if $A$ is
surjective, then $R_0=R_1=X$, so the descent is zero.  Equal ascent and descent
then gives $N_0=N_1$, hence $\Ker A=\{0\}$.  A bounded bijection between
Banach spaces has bounded inverse.  Applying the same reasoning to the compact
operator $T^*$ gives the adjoint statement.
\end{proof}

\begin{notelemma}[Biorthogonal system]
If $f_1,\ldots,f_m\in X^*$ are linearly independent, there are
$z_1,\ldots,z_m\in X$ such that
\[
  f_j(z_k)=\delta_{jk}.
\]
\end{notelemma}

\begin{proof}
Proceed by induction.  Suppose $z_1,\ldots,z_m$ have been constructed.  The
functional $f_{m+1}$ cannot vanish on
$\bigcap_{j=1}^{m}\Ker f_j$, for otherwise it would be a linear combination
of $f_1,\ldots,f_m$.  Choose $w$ in that intersection with
$f_{m+1}(w)\ne0$ and put $z_{m+1}=w/f_{m+1}(w)$.  Replacing each old $z_j$ by
$z_j-f_{m+1}(z_j)z_{m+1}$ preserves the previous relations and makes the new
one biorthogonal as well.
\end{proof}

\begin{notetheorem}[Equality of nullities]
For $A=T-\lambda I$ with $\lambda\ne0$,
\[
  \dim\Ker A=\dim\Ker A^*<\infty.
\]
Equivalently, $A$ has Fredholm index zero.
\end{notetheorem}

\begin{proof}
Let $r$ be the common ascent and descent.  The Riesz decomposition gives
\[
  X=N\oplus R,
  \qquad N=\Ker A^r,
  \quad R=\Ran A^r.
\]
The restriction of $A$ to $R$ is bijective, while $N$ is finite-dimensional
and $A|_N$ is nilpotent.  Therefore finite-dimensional rank--nullity gives
\[
  \operatorname{codim}\Ran A=\dim\Ker A.
\]
Because $\Ran A$ is closed,
$\Ker A^*=(\Ran A)^\perp$, whose dimension equals
$\dim(X/\Ran A)$.  This proves the equality.
\end{proof}

\begin{notecorollary}
For a compact operator on a complex Banach space and $\lambda\ne0$,
\[
  \lambda\in\sigma_p(T)
  \quad\Longleftrightarrow\quad
  \lambda\in\sigma_p(T^*).
\]
In particular, every nonzero spectral value of $T$ is an eigenvalue.
\end{notecorollary}

\section{The Fredholm alternative and integral equations}

\begin{notedefinition}[Fredholm alternative]
A bounded operator $B\colon X\to X$ is said to satisfy the Fredholm
alternative if exactly one of the following cases occurs:
\begin{enumerate}
  \item $B$ and $B^*$ are bijective, so $Bx=y$ and $B^*f=g$ have
        unique solutions for every $y\in X$ and $g\in X^*$;
  \item the two homogeneous equations have finite-dimensional nonzero
        solution spaces of the same dimension, and the nonhomogeneous
        equations are solvable precisely under the orthogonality conditions
        \[
          f(y)=0\quad(f\in\Ker B^*),
          \qquad
          g(x)=0\quad(x\in\Ker B).
        \]
\end{enumerate}
\end{notedefinition}

\begin{notetheorem}[Fredholm alternative for compact operators]
If $T\colon X\to X$ is compact and $\lambda\ne0$, then
$T-\lambda I$ satisfies the Fredholm alternative.
\end{notetheorem}

\begin{proof}
This is the combination of closed range, the two solvability criteria,
injectivity--surjectivity equivalence, and equality of nullities established
above.
\end{proof}

\begin{notetheorem}[Arzel\`a--Ascoli, sequential form]
A uniformly bounded and equicontinuous sequence in $C([a,b])$ has a uniformly
convergent subsequence.
\end{notetheorem}

\begin{notetheorem}[Compact integral operator]
Let $k\in C([a,b]\times[a,b])$ and define
\[
  (Tx)(s)=\int_a^b k(s,t)x(t)\,\mathrm dt
  \qquad(x\in C([a,b])).
\]
Then $T\in\Kop{C([a,b])}{C([a,b])}$.
\end{notetheorem}

\begin{proof}
Linearity is immediate, and
\[
  \|Tx\|_\infty
  \le \|x\|_\infty
      \max_{s\in[a,b]}\int_a^b|k(s,t)|\,\mathrm dt.
\]
Let $(x_n)$ be bounded by $M$ and put $y_n=Tx_n$.  Uniform continuity of $k$
on the compact square implies that, whenever $|s_1-s_2|$ is sufficiently
small,
\[
  |y_n(s_1)-y_n(s_2)|
  \le M\int_a^b|k(s_1,t)-k(s_2,t)|\,\mathrm dt
  <\varepsilon
\]
uniformly in $n$.  Thus $(y_n)$ is uniformly bounded and equicontinuous, so
Arzel\`a--Ascoli provides a uniformly convergent subsequence.
\end{proof}

\begin{notecorollary}[Fredholm integral equation]
For $\mu\ne0$, consider
\[
  x(s)-\mu\int_a^b k(s,t)x(t)\,\mathrm dt=\widetilde y(s).
\]
Exactly one of the following occurs:
\begin{enumerate}
  \item for every $\widetilde y\in C([a,b])$ there is a unique solution;
  \item the homogeneous equation has a finite-dimensional nonzero solution
        space, and the nonhomogeneous equation is solvable exactly when its
        right-hand side annihilates every solution of the adjoint homogeneous
        equation in the dual space.
\end{enumerate}
On $L^2([a,b])$, when $k\in L^2([a,b]^2)$, the same conclusion follows from
compactness of the Hilbert--Schmidt operator; its Hilbert-space adjoint has
kernel $\overline{k(t,s)}$.
\end{notecorollary}



\end{document}
